\documentclass[10pt]{amsart}

\usepackage[T1]{fontenc}
\usepackage{lmodern}
\usepackage[protrusion=true,expansion=false]{microtype}
\usepackage{mathtools}
\usepackage{amssymb}
\usepackage{mathrsfs}
\usepackage{booktabs}
\usepackage{enumitem}
\usepackage{xcolor}
\usepackage[colorlinks=true,linkcolor=blue!45!black,citecolor=blue!45!black,urlcolor=blue!55!black]{hyperref}
\usepackage[nameinlink,capitalize,noabbrev]{cleveref}

\numberwithin{equation}{section}

\newtheorem{theorem}{Theorem}[section]
\newtheorem{proposition}[theorem]{Proposition}
\newtheorem{lemma}[theorem]{Lemma}
\newtheorem{corollary}[theorem]{Corollary}
\theoremstyle{definition}
\newtheorem{definition}[theorem]{Definition}
\theoremstyle{remark}
\newtheorem{remark}[theorem]{Remark}

\newcommand{\CC}{\mathbf C}
\newcommand{\QQ}{\mathbf Q}
\newcommand{\PP}{\mathbf P}
\newcommand{\GGm}{\mathbf G_m}
\newcommand{\IC}{\operatorname{IC}}
\newcommand{\Nm}{\operatorname{Nm}}
\newcommand{\Alt}{\operatorname{Alt}}
\newcommand{\Sym}{\operatorname{Sym}}
\newcommand{\Stab}{\operatorname{Stab}}
\newcommand{\Sing}{\operatorname{Sing}}
\newcommand{\id}{\operatorname{id}}
\newcommand{\reg}{\mathrm{reg}}
\newcommand{\exc}{\mathrm{exc}}
\newcommand{\cH}{\mathcal H}
\newcommand{\cI}{\mathcal I}
\newcommand{\cM}{\mathcal M}
\newcommand{\cN}{\mathcal N}
\newcommand{\cB}{\mathcal B}
\newcommand{\cP}{\mathcal P}
\newcommand{\cX}{\mathcal X}
\newcommand{\cR}{\mathcal R}
\newcommand{\cU}{\mathcal U}
\newcommand{\cA}{\mathcal A}
\newcommand{\cC}{\mathcal C}
\newcommand{\cE}{\mathcal E}
\newcommand{\cL}{\mathcal L}
\newcommand{\wh}{\widehat}
\newcommand{\wt}{\widetilde}
\newcommand{\bo}{\boxtimes}

\title[A realization of $G_2$ for Litt Problem No. 13]
      {A Realization of $G_2$ for Litt Problem No. 13}
\author{OpenAI}
\date{Revised August 10, 2026}

\hypersetup{
  pdftitle={A Realization of G2 for Litt Problem No. 13},
  pdfauthor={OpenAI}
}

\subjclass[2020]{14K12, 14H40, 18M25, 20G41}
\keywords{abelian variety, cyclic cover, Prym variety, intersection complex,
convolution, Tannaka group, exceptional group $G_2$}

\begin{document}
\raggedbottom

\begin{abstract}
A family of complex abelian threefolds is constructed from cyclic triple covers
\[
  C:\quad w^3=y-c,\qquad c\ne0,
\]
of elliptic curves $y^2=x^3+Ax+B$.  The induced Prym polarization has type
$(1,1,3)$.  The centred theta section $X$ is a normal symmetric surface with
a unique simple elliptic singularity of type $\wt E_7$, and
$\chi(\IC_X)=7$.  For general parameters, the pair-incidence curve
$X\cap(-x-X)$ has exactly two nonexceptional components.  Their top
cohomology gives an unshifted summand $\IC_X\subset\Alt^2(\IC_X)$ in the
convolution category, while a generic unitary twist has Hodge multiplicities
$(2,3,2)$.  Together these facts determine the Tannaka pair: the connected
group is $G_2$ on its standard seven-dimensional representation, and the
same exterior-square morphism excludes the only possible disconnected
extension.  The full convolution Tannaka group is therefore $G_2$, realizing
the $G_2$ case of Litt Problem No.~13.
\end{abstract}

\maketitle

\section{Introduction}

Let $A$ be a complex abelian variety.  After negligible objects are
quotiented out, convolution turns the relevant constructible complexes on
$A$ into a neutral Tannakian category.  Thus an integral subvariety
$Z\subset A$ with $\chi(\IC_Z)\ne0$ determines a Tannaka group
$G(\IC_Z)$ and a faithful module of dimension $\chi(\IC_Z)$; see
\cite[Proposition~3.1 and Definition~3.2]{JKLM2025}.  Litt Problem No.~13 asks whether the four
exceptional simple types
\[
  G_2,\qquad E_7,\qquad E_8,\qquad F_4
\]
can occur in this manner \cite{Litt13}.  Only the $G_2$ case is considered
here; no claim is made about the other three.  The construction necessarily
uses a singular theta surface, so the exclusion results for smooth
subvarieties in \cite{KLM2026} do not apply.

The construction is defined over a rational parameter space.  Put
\[
 \Delta_E=4A^3+27B^2,
 \qquad
 \Delta_{\mathrm{br}}=4A^3+27(B-c^2)^2,
\]
and let
\[
 \cB_{\QQ}=\operatorname{Spec}
 \QQ[A,B,c,c^{-1},\Delta_E^{-1},\Delta_{\mathrm{br}}^{-1}],
 \qquad
 \cB=\cB_{\QQ}\times_{\QQ}\CC.
\]
For $b=(A,B,c)\in\cB(\CC)$, let
\[
 E_b:\ y^2=x^3+Ax+B,
 \qquad
 C_b:\ w^3=y-c.
\]
Here $C_b$ denotes the smooth projective cyclic cover defined by the section
$y-c\in H^0(E_b,\mathcal O_{E_b}(3O))$; the displayed equation is its
affine chart.  The reduced branch divisor makes this cover connected.  Let
\[
 \pi_b:C_b\longrightarrow E_b,
 \qquad (x,y,w)\longmapsto(x,y)
\]
be the cyclic triple cover.  The condition $\Delta_{\mathrm{br}}\ne0$
means that the line $y=c$ meets $E_b$ in three distinct points; these are
precisely the branch points of $\pi_b$.  Write $O\in E_b$ for the point at
infinity and set
\[
 \kappa_b=\pi_b^*\mathcal O_{E_b}(O),
 \qquad
 P_b=\ker(\Nm_{\pi_b})^0.
\]
The centred theta divisor on $J(C_b)$ restricts to the divisor
\[
 X_b=\{M\in P_b:h^0(C_b,\kappa_b\otimes M)>0\}.
\]
Write $\cP\to\cB$ and $\cX\subset\cP$ for the resulting relative Prym and
relative theta section.

\begin{theorem}[Main theorem]\label{thm:main}
There is a nonempty Zariski-open subset $\cB^\circ\subset\cB$ such that for
every $b\in\cB^\circ(\CC)$ the following hold.
\begin{enumerate}[label=\textup{(\roman*)},leftmargin=*]
\item $P_b$ is an abelian threefold, and the polarization induced by the
      centred theta divisor has type $(1,1,3)$.
\item $X_b\subset P_b$ is an integral normal symmetric ample Cartier divisor
      with a unique singular point at the origin; this singularity is simple
      elliptic of type $\wt E_7$.
\item $\chi(\IC_{X_b})=7$, and the associated Tannaka pair satisfies
\[
  \bigl(G(\IC_{X_b}),\omega(\IC_{X_b})\bigr)\simeq(G_2,V_7),
\]
      where $V_7$ is the standard seven-dimensional representation of
      $G_2$.  In particular, the full convolution Tannaka group is connected
      and isomorphic to $G_2$.
\end{enumerate}
Moreover, $\cX\subset\cP$ is a symmetric relatively ample effective Cartier
divisor, flat over $\cB$, with geometrically integral fibres.
\end{theorem}

\paragraph{Proof outline.}
Fix $b\in\cB(\CC)$ and write $(P,X)=(P_b,X_b)$.  The proof has a geometric
part and a Tannakian part, with the pair-incidence calculation linking the
two.  The norm fibre in $C^{(3)}$ is birational to $X$ and has one
$\wt E_6$ singularity.  Blowing up that point and contracting the strict
transform of the trigonal pencil gives the minimal resolution
$\rho:S\to X$; its exceptional locus is an elliptic curve of
self-intersection $-2$.  Adjunction, the local geometric genus, and Noether's
formula give
\[
 K_S^2=16,\qquad \chi(\mathcal O_S)=2,\qquad e(S)=8,
\]
while the theta group yields $p_g(S)=5$ and $q(S)=4$.  The semismall
decomposition then gives $\chi(\IC_X)=7$, and the unique singular point
forces $\Stab_P(X)=\{0\}$.  A generic unitary twist of the same resolution
has Hodge multiplicities $(2,3,2)$.

The geometric core of the argument is the generic pair-incidence fibre.
For a geometric generic point $x$ of
$\cM=(\cX/\cB)_{\mathrm{sm}}$, let $D_x$ be its unique effective divisor of
degree three.  The cyclic norm of a section of $H^0(C,3\kappa-D_x)$ is the
fixed norm line times a conic.  Its determinant has a structural linear
factor and a geometrically integral residual octic.  On that octic an exact
square-class identity identifies the ordering cover of the two line factors
with the constant quadratic cover $\QQ(\sqrt{-3})/\QQ$.  Hence over $\CC$
there are two branches.  The scheme-theoretic Abel inverse converts them into
the two nonexceptional components of
\[
 C_x=\{y\in X:-x-y\in X\},
\]
and a separate boundary argument shows that no further nonexceptional
component occurs.

After spreading these two components over a nonempty open
$\cM^\circ\subset\cM$, top cohomology of the pair fibres and the
decomposition theorem produce an unshifted summand
$\IC_X\subset\Alt^2(\IC_X)$.  If $G=G(\IC_X)$ and
$V=\omega(\IC_X)$, the final input is therefore
\[
 \begin{gathered}
 G\subset \mathrm{O}(V),\qquad \dim V=7,\qquad V|_{G^0}\ \text{irreducible},\\
 \mu:\GGm\longrightarrow [G^0,G^0]
 \quad\text{with weights }-2,0,2\text{ of multiplicities }2,3,2,\\
 \operatorname{Hom}_G(V,\Lambda^2V)\ne0.
 \end{gathered}
\]
The exterior-square morphism rules out $\mathrm{SO}_7$, the Hodge cocharacter rules
out the seven-dimensional $\mathrm{PGL}_2$ module, and the same morphism excludes the
only possible disconnected extension of $G_2$.  Thus $G=G_2$.

\paragraph{Organization.}
The proof is organized by the following dependency chain.
\begin{enumerate}[label=\textup{(\arabic*)},leftmargin=*]
\item \cref{prop:theta-resolution,lem:surface-invariants,cor:theta-ic,lem:conormal-degree}
      establish the geometry of the theta section and $\chi(\IC_X)=7$.
\item \cref{prop:resolution-irregularity,prop:hodge-vector} compute the
      invariants of the resolution and the generic-twist Hodge
      multiplicities $(2,3,2)$.
\item \cref{prop:determinant-factorization,lem:residual-octic-integral,%
      prop:ordering-cover,prop:abel-incidence,lem:boundary-exhaustion}
      identify the two generic pair components, and
      \cref{thm:uniform-pair-splitting} spreads them.
\item \cref{lem:split-fibre-top,lem:off-theta-vanishing,lem:extract-ic}
      extract an unshifted $\IC_X$ summand in alternating pair convolution.
\item \cref{prop:connected-g2,cor:full-g2} identify the connected and full
      Tannaka groups.
\end{enumerate}

Unless a rational model or a finite-field reduction is explicitly mentioned,
the geometric construction is over $\CC$.  Intersection complexes,
cohomology, and Hodge modules retain their canonical $\QQ$-structures.  The
Tannakian category is obtained after extension of the underlying perverse
sheaves from $\QQ$ to $\CC$.  Thus all Tannaka groups and tautological
modules below are over $\CC$.
The convention $\PP(\mathcal E)=\operatorname{Proj}(\operatorname{Sym}\mathcal E^\vee)$ is used,
so $\PP(\mathcal E)$ parametrizes lines in the fibres of $\mathcal E$.  For
an integral scheme $Y$, $k(Y)$ denotes its function field and $k(y)$ the
residue field of a point $y$.

\section{Tannakian conventions}

Let $A$ be a complex abelian variety and let $m:A\times A\to A$ be addition.
For constructible complexes with coefficients in $\QQ$ or $\CC$, set
\[
K*L=Rm_*(K\bo L).
\]
Set
\[
 \begin{aligned}
 \mathsf S(A)&=\{Q\in\operatorname{Perv}(A,\CC):\chi(A,Q)=0\},\\
 \mathsf T(A)&=\{M\in D_c^b(A,\CC):{}^pH^i(M)\in\mathsf S(A)
                   \text{ for every }i\}.
 \end{aligned}
\]
The first is a Serre subcategory and the second is the corresponding
triangulated tensor ideal.  The quotient
\[
 \overline D(A)=D_c^b(A,\CC)/\mathsf T(A)
\]
inherits the perverse $t$-structure; its heart is
\[
 \overline{\operatorname{Perv}}(A)
   =\operatorname{Perv}(A,\CC)/\mathsf S(A).
\]
Convolution is $t$-exact on $\overline D(A)$, and the displayed heart is a
rigid symmetric monoidal neutral Tannakian category
\cite[Proposition~3.1]{JKLM2025}.  For the pure polarizable objects used
below, the tensor subcategory generated by one object is semisimple by the
decomposition theorem; its Tannaka group is therefore reductive.
Let $\delta_0$ denote the skyscraper perverse sheaf at the origin; its
image is the tensor unit for convolution.

For an integral subvariety $Z\subset A$ of dimension $d$, let
\[
 \IC_Z=j_{!*}\QQ_{Z_{\reg}}[d],
 \qquad j:Z_{\reg}\hookrightarrow A,
\]
be its perverse intersection complex, extended by zero to $A$.  If
$\chi(\IC_Z)>0$, let $G(\IC_Z)$ denote the Tannaka group of the category
generated by $\IC_Z\otimes_{\QQ}\CC$ in the displayed quotient, and by
$\omega(\IC_Z)$ its tautological module.  Then
\[
 \dim\omega(\IC_Z)=\chi(\IC_Z).
\]
The action on the tautological module is faithful by construction.  Scalar
extension is suppressed inside Tannakian expressions.  Let $\IC_Z^H$ denote
the polarizable pure Hodge module whose underlying
perverse $\QQ$-sheaf is $\IC_Z$.

The commutativity constraint contains the usual cohomological sign.  Since
the support in this paper has even dimension two, the alternating projector
on $\IC_X*\IC_X$ is carried by the fibre functor to the ordinary exterior
square of $\omega(\IC_X)$.  Let $\Alt^2(\IC_X)$ denote its image in the
Tannakian quotient.

\section{The shifted cyclic theta section}

Until \cref{sec:spreading}, fix $b=(A,B,c)\in\cB(\CC)$ and suppress it from
notation.  Thus
\[
 E:y^2=f_3(x),\qquad C:w^3=y-c,
\]
where $f_3(x)=x^3+Ax+B$.  Set
\[
 Y=y-c,
\]
so that $w^3=Y$.  The map $\pi:C\to E$ is a cyclic cover of degree three.
Its branch divisor is the reduced intersection of $E$ with the line $y=c$.

\subsection{The cover, the Prym, and the polarization}

Let $b_1,b_2,b_3\in E$ be the branch points and let $r_i\in C$ be the
ramification point above $b_i$.  Since
\[
 \operatorname{div}_E(y-c)=b_1+b_2+b_3-3O,
\]
the branch divisor belongs to $|3O|$ and, in the group law of $E$,
\begin{equation}
 b_1+b_2+b_3=O.
\label{eq:branch-sum}
\end{equation}
Put $\eta=\mathcal O_E(O)$.  The cyclic covering algebra is
\[
 \pi_*\mathcal O_C
 =\mathcal O_E\oplus\eta^{-1}\oplus\eta^{-2}.
\]
Riemann--Hurwitz gives $g(C)=4$, and the canonical bundle formula gives
\[
 K_C=\pi^*(K_E\otimes\eta^2)=\pi^*\eta^2=\kappa^2,
 \qquad \kappa=\pi^*\eta.
\]
By projection formula,
\[
 h^0(C,\kappa)
 =h^0(E,\eta)+h^0(E,\mathcal O_E)+h^0(E,\eta^{-1})=2.
\]

The pullback $\pi^*:J(E)\to J(C)$ is injective.  Indeed, if
$\alpha\in\operatorname{Pic}^0(E)$ and $\pi^*\alpha\simeq\mathcal O_C$, then
\[
 1=h^0(C,\pi^*\alpha)
  =h^0(E,\alpha)+h^0(E,\alpha\eta^{-1})
   +h^0(E,\alpha\eta^{-2}),
\]
and the last two terms vanish; hence $\alpha=\mathcal O_E$.  By duality,
$\ker(\Nm_\pi)$ is connected.  Thus set
\[
 P=\ker(\Nm_\pi).
\]
It has dimension three.  If $\tau(w)=\zeta_3w$, then
$\pi^*\Nm_\pi=1+\tau+\tau^2$ on $J(C)$.  Consequently
$\operatorname{Im}(1-\tau)\subset P$.  The $\tau$-invariant part of
$H^0(C,\Omega_C^1)$ is $\pi^*H^0(E,\Omega_E^1)$ and has dimension one, so
$\operatorname{Im}(1-\tau)$ has dimension three.  Hence
\begin{equation}
 P=\operatorname{Im}(1-\tau)=\ker(\Nm_\pi).
\label{eq:prym-conventions}
\end{equation}
The same argument works in the relative family.  Indeed
$\Nm_\pi\circ\pi^*=[3]$ on the relative elliptic Jacobian, so the norm map
is surjective; its kernel is smooth, and the fibrewise injectivity of
$\pi^*$ proved above makes every kernel fibre connected.  Thus
$\cP=\ker(\Nm_\pi)\to\cB$ is an abelian scheme of relative dimension
three.

Use the theta characteristic $\kappa$ to centre the theta divisor on
$J(C)=\operatorname{Pic}^0(C)$:
\[
 \Theta_\kappa=\{M\in J(C):h^0(C,\kappa\otimes M)>0\}.
\]
Let $L=\mathcal O_{J(C)}(\Theta_\kappa)|_P$.  The type of the restricted
polarization can be read off directly.  Under the principal polarizations of
the two Jacobians, $\pi^*$ and $\Nm_\pi$ are dual homomorphisms; hence the
orthogonal complement of $P$ in $J(C)$ is $\pi^*J(E)$
\cite[Section~12.3]{BirkenhakeLange2004}.  Therefore
\[
 \ker(\varphi_L)=P\cap\pi^*J(E).
\]
Since $\pi^*$ is injective and $\Nm_\pi\circ\pi^*=[3]$ on $J(E)$,
\[
 P\cap\pi^*J(E)=\pi^*J(E)[3].
\]
This group has order $9$.  If the polarization type is
$(d_1,d_2,d_3)$ with $d_1\mid d_2\mid d_3$, then
$|\ker(\varphi_L)|=(d_1d_2d_3)^2$; hence $d_1d_2d_3=3$ and necessarily
\begin{equation}
 \operatorname{type}(L)=(1,1,3),
 \qquad L^3=3!\,d_1d_2d_3=18.
\label{eq:prym-polarization}
\end{equation}
Scheme-theoretically,
\[
 X=\Theta_\kappa\cap P
  =\{M\in P:h^0(C,\kappa\otimes M)>0\}.
\]
The theta divisor is preserved as a closed subscheme by the Serre-duality
involution $L'\mapsto K_C\otimes(L')^{-1}$ on $\operatorname{Pic}^3(C)$.
Since $\kappa^2=K_C$, this involution is $M\mapsto M^{-1}$ in centred
coordinates.  Relative Serre duality preserves the determinantal ideal of
theta, so after restriction to $P$ one obtains
\begin{equation}
 X=-X
\label{eq:symmetry}
\end{equation}
as closed subschemes.  The centre will be seen to be unique in
\eqref{eq:stabilizer}.

\subsection{The norm fibre and the resolution}

Let
\[
 n:C^{(3)}\longrightarrow\operatorname{Pic}^3(E),
 \qquad D\longmapsto\mathcal O_E(\pi_*D),
\]
and put
\[
 N=n^{-1}(\eta^3),
 \qquad
 a:N\longrightarrow X,
 \qquad
 D\longmapsto\mathcal O_C(D)\otimes\kappa^{-1}.
\]
For the centred Abel map $D\mapsto\mathcal O_C(D)\otimes\kappa^{-1}$ one has
scheme-theoretically
\[
 N=C^{(3)}\times_{J(C)}P.
\]
Thus $a$ is its base change to $P$.  The map $a$ is surjective: if $M\in X$,
a nonzero section of $\kappa\otimes M$ has an effective divisor $D$ of
degree three, and
\[
 \Nm_\pi\mathcal O_C(D)
 =\Nm_\pi(\kappa\otimes M)=\eta^3.
\]
The finite morphism $C^{(3)}\to E^{(3)}$ induced by $\pi$ restricts to a
finite surjection
\[
 N\longrightarrow|\eta^3|\simeq\PP^2.
\]
Hence $\dim N=2$, so the restriction of the centred theta section to $P$ is
nonzero.  Therefore $X\in|L|$ is an effective Cartier divisor, and it is
ample.

Let $\Gamma=|\kappa|\simeq\PP^1$ and put
\[
 D_0=r_1+r_2+r_3\in\Gamma.
\]

\begin{lemma}[The norm fibre]\label[lemma]{lem:norm-fibre}
The surface $N$ has the unique singular point $D_0$, and the germ $(N,D_0)$
is a simple elliptic singularity of type $\wt E_6$.  If
$\beta:\widehat N\to N$ is the blowup of $D_0$, then $\widehat N$ is smooth,
its exceptional curve $E_0$ is elliptic with $E_0^2=-3$, and the strict
transform $\widetilde\Gamma$ is a $(-1)$-curve meeting $E_0$ transversely in
one point.  Contracting $\widetilde\Gamma$ gives a smooth surface $S$ and a
proper morphism
\[
 \rho:S\longrightarrow X.
\]
The image $E_{\exc}$ of $E_0$ satisfies $E_{\exc}^2=-2$ and maps to the
origin of $X$.
\end{lemma}

\begin{proof}
Because the three branch points are distinct, a neighbourhood of $D_0$ in
$C^{(3)}$ is the product of neighbourhoods of the $r_i$.  Choose additive
analytic coordinates $t_i$ on $E$ at $b_i$ so that the Abel map
$E^{(3)}\to\operatorname{Pic}^3(E)$ is, after translation, given by
$t_1+t_2+t_3$.  After rescaling local parameters $z_i$ at $r_i$, the cover
has $t_i=z_i^3$.  Hence the norm fibre has the analytic equation
\[
 z_1^3+z_2^3+z_3^3=0.
\]
Thus $(N,D_0)$ is the affine cone over a smooth plane cubic, hence a simple
elliptic singularity of type $\wt E_6$.

It remains to locate the other critical points of the norm map.  Separate
the possible multiplicity patterns on the symmetric cube.  A divisor
containing an unramified point
has a tangent direction on which $d\pi$ is nonzero, so it cannot be critical.
Near a ramification point the map is $z\mapsto z^3$.  If two points collide,
then in the elementary symmetric coordinates $e_1=z_1+z_2$ and
$e_2=z_1z_2$ one has
\[
 z_1^3+z_2^3=e_1^3-3e_1e_2,
\]
which has no linear term at the double ramification cycle.  If three points
collide, however,
\[
 z_1^3+z_2^3+z_3^3=e_1^3-3e_1e_2+3e_3,
\]
and the $3e_3$ term makes the differential nonzero.  Thus, apart from the
cycle $D_0=r_1+r_2+r_3$, the only possible critical cycles are
$2r_i+r_j$ with $i\ne j$.  Membership in the norm fibre would then force
\[
 2b_i+b_j\sim3O.
\]
If $k$ is the third index, \eqref{eq:branch-sum} gives
$b_i+b_j+b_k\sim3O$, hence $b_i=b_k$, contradicting the fact that the branch
points are distinct.  Therefore $D_0$ is the only singular point of $N$.

Blowing up the vertex of the cubic cone gives a smooth surface $\widehat N$
with exceptional elliptic curve $E_0$ and $E_0^2=-3$.  The local behaviour
of the pencil $\Gamma=|\kappa|$ at the vertex is as follows.  Projection
formula gives
\[
 \pi_*\kappa\simeq\eta\oplus\mathcal O_E\oplus\eta^{-1}.
\]
The two sections of $\kappa$ may therefore be chosen as follows: the section
coming from the $\mathcal O_E$-summand is represented by $w$ and has zero
divisor $D_0$, while a generator of $H^0(E,\eta)$ pulls back to a section
that is nonzero at each $r_i$, since none of the branch points $b_i$ is $O$.  Let $t$ be the affine parameter of the pencil
centred at the member $D_0$.  In a local trivialization of $\kappa$ at
$r_i$, the moving section has equation
\[
 w+t u_i=0,\qquad u_i(r_i)\ne0.
\]
Since $w$ is a local parameter at the totally ramified point $r_i$, after
rescaling $z_i$ this gives
\[
 z_i(t)=-u_i(r_i)t+O(t^2).
\]
Thus $\Gamma\to N$ has a nonzero tangent vector at $D_0$ and hence
multiplicity one there.  Because $\Gamma\subset N$, its tangent direction
lies on the exceptional cubic
\[
 \lambda_1^3+\lambda_2^3+\lambda_3^3=0,
 \qquad \lambda_i=-u_i(r_i),
\]
and the displayed expansion shows that the strict transform of $\Gamma$
meets $E_0$ in one reduced point, transversely.

Let $a_C:C^{(3)}\to J(C)$ be the Abel--Jacobi map and let $\Theta_C$ be the
principal theta divisor.  If $X_{p_0}\subset C^{(d)}$ denotes the divisor of
cycles containing a fixed point $p_0$, Macdonald's canonical-class formula is
\cite{Macdonald1962}
\[
 K_{C^{(d)}}\equiv a_d^*\Theta_C+(g-d-1)X_{p_0}.
\]
For $g(C)=4$ and $d=3=g-1$ the second term vanishes, hence
\[
 K_{C^{(3)}}\equiv a_C^*\Theta_C.
\]
Since $N$ is a fibre of $n:C^{(3)}\to\operatorname{Pic}^3(E)$, its normal
line bundle is trivial.  Adjunction therefore gives
\[
 K_N\equiv a^*L.
\]
The map $a$ contracts $\Gamma$, so $K_N\cdot\Gamma=0$.  The cubic cone has
discrepancy $-1$, hence
\[
 K_{\widehat N}=\beta^*K_N-E_0.
\]
Since $E_0\cdot\widetilde\Gamma=1$, one has
$K_{\widehat N}\cdot\widetilde\Gamma=-1$.  Adjunction on
$\widetilde\Gamma\simeq\PP^1$ then gives
$\widetilde\Gamma^2=-1$.  Contracting it produces a smooth surface $S$ and
\[
 E_{\exc}^2=E_0^2+(E_0\cdot\widetilde\Gamma)^2=-2.
\]
The composite $\widehat N\to N\to X$ is constant on
$\widetilde\Gamma$, hence factors through $\rho:S\to X$.  The curve
$E_{\exc}$ maps to the origin.
\end{proof}

\begin{proposition}[The theta surface and its resolution]
\label[proposition]{prop:theta-resolution}
The theta section $X\subset P$ is an integral normal symmetric ample Cartier
divisor with $\Sing(X)=\{0\}$.  The morphism $\rho:S\to X$ of
\cref{lem:norm-fibre} is its minimal resolution and is an isomorphism over
$X\setminus\{0\}$.  Its exceptional locus is the smooth elliptic curve
$E_{\exc}$ with $E_{\exc}^2=-2$; in particular, $(X,0)$ is a simple elliptic
singularity of type $\wt E_7$.
\end{proposition}

\begin{proof}
A basis of $H^0(C,K_C)$ is represented on the affine chart by
\[
 \frac{dx}{y},\qquad
 \frac{dx}{wy},\qquad
 \frac{dx}{w^2y},\qquad
 \frac{x\,dx}{w^2y}.
\]
Indeed these forms are regular both at the ramification points and at the
three points above $O$.  After multiplication by the common factor
$w^2y/dx$, the canonical map is
\[
 (x,w)\longmapsto[w^2:w:1:x].
\]
The map is birational, so $C$ is not hyperelliptic.  Among the ten
quadratic monomials in the four canonical coordinates there is the evident
relation $X_0X_2-X_1^2=0$, and it is the only one.  After restriction to
$C$, the other nine monomials may be represented as
\[
 \underbrace{1,x,Y,x^2}_{\tau\text{-eigenvalue }1};\qquad
 \underbrace{w,xw,wY}_{\tau\text{-eigenvalue }\zeta_3};\qquad
 \underbrace{w^2,xw^2}_{\tau\text{-eigenvalue }\zeta_3^2}.
\]
At any of the three points above $O$ their pole orders within these three
eigenspaces are, respectively,
\[
 (0,2,3,4),\qquad (1,3,4),\qquad (2,4).
\]
They are therefore linearly independent within each eigenspace, and terms
from distinct eigenspaces cannot cancel.  Thus the canonical ideal has a
single quadric,
\[
 X_0X_2-X_1^2=0,
\]
which has rank three and one ruling.  By geometric Riemann--Roch, every
$g^1_3$ on a nonhyperelliptic genus-four curve determines a ruling of the
canonical quadric.  Hence $\kappa$ is the unique $g^1_3$ on $C$.
Riemann--Kempf \cite{Kempf1973} shows that the Abel map is an isomorphism
over the locus $h^0=1$.  The corresponding points of $N$ are smooth by the
critical-point calculation above, and $|\kappa|$ is the only
positive-dimensional Abel fibre.  It follows that
$\Sing(X)\subset\{0\}$.

The Riemann singularity theorem \cite{Kempf1973} gives
$\operatorname{mult}_0(\Theta_\kappa)=h^0(C,\kappa)=2$.  The restricted
theta equation is nonzero, since $\dim X=2<\dim P$, and lies in
$\mathfrak m_{P,0}^2$.  Thus $X$ is singular at the origin, and
\begin{equation}
 \Sing(X)=\{0\}.
\label{eq:singular-locus}
\end{equation}
The Cartier divisor $X$ is generically reduced and Cohen--Macaulay, hence
reduced.  Since $L$ is ample, $X$ is connected.  Suppose that $X$ were
reducible.  Two irreducible components joined by an edge in the incidence
graph must then meet.  Both have dimension two in the smooth threefold $P$,
so every nonempty intersection of two distinct components has dimension at
least one.  At such a point the local ring of $X$ has at least two minimal
primes and is therefore not regular.  Hence the intersection is contained in
$\Sing(X)=\{0\}$, contradicting its positive dimension.  Thus $X$ is
integral.  Being Cohen--Macaulay, it satisfies $S_2$, while
\eqref{eq:singular-locus} gives regularity in codimension one.  Serre's
criterion makes $X$ normal.

The uniqueness of the $g^1_3$ shows that the Abel map is an isomorphism away
from $\Gamma$.  Hence $\rho$ is a resolution whose only exceptional curve is
$E_{\exc}$.  Since this curve is elliptic with self-intersection $-2$, the
resolution is minimal and the simple-elliptic classification identifies the
germ with type $\wt E_7$ \cite{Laufer1977,Pinkham1977}.
\end{proof}

\begin{lemma}[Surface invariants]\label[lemma]{lem:surface-invariants}
The minimal resolution $S$ satisfies
\[
 K_S^2=16,\qquad \chi(\mathcal O_S)=2,\qquad e(S)=8.
\]
\end{lemma}

\begin{proof}
From $K_N\equiv a^*L$ and $X\in|L|$ one obtains
\[
 K_N^2=(a^*L)^2=L^2\cdot[X]=L^3=18.
\]
For the blowup of the $\wt E_6$ vertex,
$K_{\widehat N}=\beta^*K_N-E_0$, so
\[
 K_{\widehat N}^2=18+E_0^2=15.
\]
The contraction $\widehat N\to S$ is the inverse of a blowup at a smooth
point, and therefore $K_S^2=16$.

The Cartier sequence
\[
 0\longrightarrow L^{-1}\longrightarrow\mathcal O_P
 \longrightarrow\mathcal O_X\longrightarrow0
\]
gives $\chi(\mathcal O_X)=3$: on the abelian threefold $P$,
$\chi(\mathcal O_P)=0$, $\chi(L)=3$, and
$\chi(L^{-1})=-\chi(L)=-3$.  A simple elliptic singularity has local
geometric genus one; equivalently,
$R^1\rho_*\mathcal O_S$ is a skyscraper sheaf of length one
\cite{Laufer1977}.  Since $X$ is normal,
$\rho_*\mathcal O_S=\mathcal O_X$, and Leray gives
\[
 \chi(\mathcal O_S)=\chi(\mathcal O_X)-1=2.
\]
Noether's formula now yields
\[
 e(S)=12\chi(\mathcal O_S)-K_S^2=24-16=8.
\]
\end{proof}

\begin{corollary}[The intersection complex]
\label[corollary]{cor:theta-ic}
There is a decomposition
\[
 R\rho_*\QQ_S[2]\simeq\IC_X\oplus\delta_0.
\]
Moreover,
\[
 \chi(\IC_X)=7,
 \qquad
 \Stab_P(X)=\{0\}.
\]
\end{corollary}

\begin{proof}
The map $\rho:S\to X$ is semismall, with the single relevant exceptional
fibre $E_{\exc}$.  Its top intersection form is the nondegenerate matrix
$(-2)$.  The semismall decomposition theorem
\cite[Theorem~3.4.1]{deCataldoMigliorini2002} gives
\begin{equation}
 R\rho_*\QQ_S[2]\simeq\IC_X\oplus\delta_0.
\label{eq:semismall}
\end{equation}
Taking Euler characteristics and using \cref{lem:surface-invariants} gives
\begin{equation}
 \chi(\IC_X)=e(S)-1=7.
\label{eq:chi-seven}
\end{equation}
A translation preserving $X$ must preserve its unique singular point, hence
\begin{equation}
 \Stab_P(X)=\{0\}.
\label{eq:stabilizer}
\end{equation}
In particular, the origin is the unique centre of the symmetry
\eqref{eq:symmetry}.
\end{proof}

Let $\Lambda_X\subset T^*P$ be the closure of the conormal bundle of
$X_{\reg}$, and let $\Lambda_0=T_0^*P$ be the conormal fibre at the origin.
Let
\[
 \gamma_X:\PP(\Lambda_X)\longrightarrow\PP(T_0^*P)
\]
be the map induced by cotangent projection.

\begin{lemma}[Conormal degree]\label[lemma]{lem:conormal-degree}
The characteristic cycle and conormal Gauss degree are
\[
 \operatorname{CC}(\IC_X)=\Lambda_X+\Lambda_0,
 \qquad
 \deg(\gamma_X)=6.
\]
In particular, the conormal Gauss map is dominant.
\end{lemma}

\begin{proof}
The sectional calculation follows directly from the
standard analytic form of a $\wt E_7$ simple-elliptic germ,
\[
 z^2+f_4(x,y)=0,
\]
where $f_4$ is a square-free binary quartic
\cite{Laufer1977,Pinkham1977}.  On a general hyperplane $z=ax+by$, the
quadratic term has rank one and the restriction of $f_4$ to its null line is
nonzero.  The analytic splitting lemma therefore reduces the plane-curve
germ to $s^2+t^4=0$, namely type $A_3$, whose Milnor number is $3$.
The Lefschetz formula for the local Euler obstruction
\cite[Theorem~3.1]{BrasseletLeSeade2000} identifies $\operatorname{Eu}_X(0)$
with the Euler characteristic of this generic linear Milnor fibre.  Hence
\[
 \operatorname{Eu}_X(0)=1-\mu(A_3)=1-3=-2.
\]
On the other hand, the stalk of $R\rho_*\QQ_S[2]$ at the origin has
Euler characteristic $e(E_{\exc})=0$, whereas the summand $\delta_0$ has
stalk Euler characteristic $1$.  Thus \eqref{eq:semismall} gives
$\chi((\IC_X)_0)=-1$.  The Dubson--Kashiwara local index formula
\cite{Kashiwara1985} therefore gives
\[
 \operatorname{CC}(\IC_X)=\Lambda_X+m\Lambda_0,
 \qquad -1=-2+m,
\]
hence $m=1$.  The Gauss map of $\Lambda_0=T_0^*P$ is the identity on
$\PP(T_0^*P)$ and has degree one.  The Franecki--Kapranov index theorem
\cite[Theorem~1.3]{FraneckiKapranov2000}, together with
\eqref{eq:chi-seven}, then gives
\[
 7=\deg(\gamma_X)+1.
\]
Thus
\begin{equation}
 \deg(\gamma_X)=6.
\label{eq:gauss-degree}
\end{equation}
\end{proof}

\subsection{The irregularity and the generic-twist Hodge vector}

The exceptional elliptic curve contributes to $H^1(S)$.  We compute this
contribution before considering twists.

\begin{proposition}[Irregularity of the resolution]
\label[proposition]{prop:resolution-irregularity}
The minimal resolution $S$ satisfies
\[
 p_g(S)=5,\qquad q(S)=4.
\]
Consequently
\[
 b_1(S)=8,\qquad b_2(S)=22,\qquad h^{1,1}(S)=12.
\]
\end{proposition}

\begin{proof}
Recall from \eqref{eq:prym-polarization} that
\[
 K(L):=\ker(\varphi_L)=\pi^*E[3]
\]
has order nine.  Every point of $K(L)$ lies on $X$.  For
$\alpha\in E[3]$, the projection formula and
$\pi_*\mathcal O_C=\mathcal O_E\oplus\eta^{-1}\oplus\eta^{-2}$ give
\[
\begin{aligned}
 h^0(C,\kappa\otimes\pi^*\alpha)
 &=h^0(E,\eta\otimes\alpha)+h^0(E,\alpha)
   +h^0(E,\eta^{-1}\otimes\alpha).
\end{aligned}
\]
The first term is one because $\deg\eta=1$; the other two vanish when
$\alpha\ne\mathcal O_E$.  For $\alpha=\mathcal O_E$ the same expression is
two.  Thus
\begin{equation}
 K(L)\subset X.
\label{eq:kernel-on-theta}
\end{equation}

Let $s\in H^0(P,L)$ be a defining section of $X$.  The theta group of $L$
acts on the three-dimensional space $H^0(P,L)$ through its irreducible
Schr\"odinger representation; see, for example,
\cite[Chapter~6]{BirkenhakeLange2004}.  For $a\in K(L)$ choose a lift of
translation by $a$ to the theta group.  The corresponding translate of $s$
vanishes at the origin, because its value there is a nonzero scalar multiple
of $s(a)$ and \eqref{eq:kernel-on-theta} gives $s(a)=0$.  The orbit of the
nonzero vector $s$ spans $H^0(P,L)$ by irreducibility.  Hence
\begin{equation}
 t(0)=0\qquad\text{for every }t\in H^0(P,L).
\label{eq:basepoint-origin}
\end{equation}
In other words, the origin is a base point of the complete linear system
$|L|$.

Adjunction gives $\omega_X\simeq L|_X$, and the divisor sequence is
\[
 0\longrightarrow\mathcal O_P\xrightarrow{\,s\,}L
 \longrightarrow\omega_X\longrightarrow0.
\]
Since $L$ is ample on the abelian threefold, $H^i(P,L)=0$ for $i>0$.
Therefore
\begin{equation}
 0\longrightarrow \CC s\longrightarrow H^0(P,L)
 \longrightarrow H^0(X,\omega_X)
 \longrightarrow H^1(P,\mathcal O_P)\longrightarrow0,
\label{eq:canonical-sections-exact}
\end{equation}
so $h^0(X,\omega_X)=2+3=5$.

All five canonical sections vanish at the singular point.
The two-dimensional subspace coming from $H^0(P,L)/\CC s$ does so by
\eqref{eq:basepoint-origin}.  For the remaining three directions, let $v$ be a translation-invariant
vector field on $P$.  Choose local frames of $L$ with transition functions
$g_{ij}$ and write the defining section as local functions $s_i$, so
$s_i=g_{ij}s_j$.  Differentiating gives
\[
 v(s_i)-g_{ij}v(s_j)=v(g_{ij})s_j.
\]
The right-hand side vanishes on $X$, so the restrictions of the $v(s_i)$
glue to a well-defined section
$\partial_v s\in H^0(X,L|_X)$.  Under the boundary map in
\eqref{eq:canonical-sections-exact}, this section maps, up to the usual sign
convention, to
\[
 \iota_v c_1(L)\in H^1(P,\mathcal O_P),
\]
because the associated \v{C}ech cocycle is $v(g_{ij})/g_{ij}$.  This is
exactly the differential at the origin of the polarization morphism
$\varphi_L:P\to\operatorname{Pic}^0(P)$.  Since $\varphi_L$ is an isogeny,
\[
 H^0(P,T_P)\longrightarrow H^1(P,\mathcal O_P),
 \qquad v\longmapsto\iota_vc_1(L),
\]
is an isomorphism.  The three sections $\partial_v s$ therefore map onto the
last term of \eqref{eq:canonical-sections-exact}.  They also vanish at the
origin because $0$ is a singular point of $X$, so $ds(0)=0$.  Thus every
section of $\omega_X$ vanishes at $0$.

Since $K_X$ is Cartier, write
$K_S=\rho^*K_X+aE_{\exc}$.  Adjunction on the exceptional elliptic curve gives
$(K_S+E_{\exc})\cdot E_{\exc}=0$, while
$\rho^*K_X\cdot E_{\exc}=0$ and $E_{\exc}^2=-2$.  Thus $a=-1$ and
\begin{equation}
 K_S=\rho^*K_X-E_{\exc}.
\label{eq:canonical-resolution}
\end{equation}
The ideal appearing after resolution is also explicit.  Since
$X$ is normal and $\rho$ is a resolution,
$\rho_*\mathcal O_S=\mathcal O_X$.  Push forward
\[
 0\longrightarrow\mathcal O_S(-E_{\exc})\longrightarrow\mathcal O_S
 \longrightarrow\mathcal O_{E_{\exc}}\longrightarrow0.
\]
The exceptional curve is connected and projective, so its global regular
functions are constant.  Since $\rho(E_{\exc})=\{0\}$, the induced map
\[
 \mathcal O_X\longrightarrow\rho_*\mathcal O_{E_{\exc}}
\]
is evaluation at the reduced point $0$ and has kernel $\mathfrak m_0$.
Left exactness of $\rho_*$ gives
\[
 \rho_*\mathcal O_S(-E_{\exc})=\mathfrak m_0.
\]
By the projection formula,
\[
 H^0(S,K_S)=H^0(X,\omega_X\otimes\mathfrak m_0).
\]
The preceding paragraph shows that the right-hand side is all of
$H^0(X,\omega_X)$, and therefore $p_g(S)=5$.  Finally
\cref{lem:surface-invariants} gives $\chi(\mathcal O_S)=2$, so
\[
 q(S)=1+p_g(S)-\chi(\mathcal O_S)=4.
\]
The stated Betti and Hodge numbers now follow from $e(S)=8$.
\end{proof}

Let $j:S\xrightarrow{\rho}X\hookrightarrow P$.  A unitary rank-one character
of $\pi_1(P)$ determines a unitary local system $\mathcal L_\alpha$ and a
holomorphic line bundle $\alpha\in\operatorname{Pic}^0(P)$.  The same symbol $\alpha$ will also denote the latter.

For $\alpha\ne\mathcal O_P$, the standard vanishing for a nontrivial
line bundle in $\operatorname{Pic}^0(P)$ and the divisor sequence give
\[
 H^0(X,(L\otimes\alpha)|_X)\simeq H^0(P,L\otimes\alpha),
 \qquad h^0(P,L\otimes\alpha)=3.
\]
Combining this with \eqref{eq:canonical-resolution} and
$\rho_*\mathcal O_S(-E_{\exc})=\mathfrak m_0$ yields
\begin{equation}
 H^0(S,K_S\otimes j^*\alpha)
 \simeq
 \ker\!\left(
 H^0(P,L\otimes\alpha)\xrightarrow{\operatorname{ev}_0}
 (L\otimes\alpha)_0\right).
\label{eq:twisted-canonical-evaluation}
\end{equation}
The locus on which the evaluation map in
\eqref{eq:twisted-canonical-evaluation} is nonzero is a nonempty Zariski-open
subset of $\operatorname{Pic}^0(P)$.  For nonemptiness, use the
surjectivity of $\varphi_L$: every $L\otimes\alpha$ is a translate of $L$,
and a translate of the defining section by a point outside $X$ does not
vanish at the origin.  Intersecting this open with its inverse under
$\alpha\mapsto\alpha^{-1}$ gives a dense open $W\subset\operatorname{Pic}^0(P)$
on which
\begin{equation}
 h^0(S,K_S\otimes j^*\alpha)
 =h^0(S,K_S\otimes j^*\alpha^{-1})=2.
\label{eq:twisted-outer-hodge}
\end{equation}

Perverse generic vanishing gives a nonempty Zariski-open subset of the Betti
character torus on which
\[
 H^k(P,Rj_*\QQ_S[2]\otimes\mathcal L_\alpha)=0\qquad(k\ne0).
\]
The compact unitary character torus is Zariski dense in the Betti character
torus, and its unitary-character identification with
$\operatorname{Pic}^0(P)$ is a real-analytic isomorphism.  The restriction of
each of the two nonempty Zariski-open loci to the unitary torus is therefore
a dense Euclidean open subset: a proper complex algebraic subset cannot
contain a Euclidean open subset of this totally real Zariski-dense torus.
Their intersection is consequently a nonempty Euclidean open subset
$\Omega$ of the unitary character torus.

\begin{proposition}[Generic-twist Hodge vector]\label[proposition]{prop:hodge-vector}
For every $\alpha\in\Omega$, the Hodge structure
$H^0(P,\IC_X^H\otimes\mathcal L_\alpha)$ has Hodge numbers
\[
 (h^{2,0},h^{1,1},h^{0,2})=(2,3,2).
\]
Equivalently, the associated antidiagonal Hodge cocharacter has weights
$-2,0,2$ on $\omega(\IC_X)$ with multiplicities $2,3,2$, respectively.
\end{proposition}

\begin{proof}
For $\alpha\in\Omega$, perverse generic vanishing
\cite[Theorem~1.1]{KraemerWeissauer2015} gives
\[
 H^k(S,j^*\mathcal L_\alpha)=0\qquad(k\ne2).
\]
For a unitary rank-one local system, Hodge theory for the associated flat
unitary line bundle gives the twisted decomposition
\[
 H^{p,q}(S,j^*\mathcal L_\alpha)
 =H^q(S,\Omega_S^p\otimes j^*\alpha),
 \qquad p+q=2.
\]
Equation \eqref{eq:twisted-outer-hodge}, together with Serre duality, gives
outer Hodge numbers $2$ and $2$.  Since the Euler characteristic of a
rank-one local system equals $e(S)=8$ and all cohomology is concentrated in
degree two, the middle Hodge number is $4$.  Thus the Hodge vector of
$Rj_*\QQ_S^H[2]$ is $(2,4,2)$.  The Hodge-module refinement of
\eqref{eq:semismall}, extended to $P$, is
\[
 Rj_*\QQ_S^H[2]\simeq\IC_X^H\oplus\QQ_0^H(-1).
\]
The point summand has type $(1,1)$, so the Hodge vector for $\IC_X^H$ is
$(2,3,2)$.
\end{proof}

The symmetry of $X$ also gives the self-duality needed in the tensor
argument.  Indeed, Verdier self-duality of the intersection complex and
\eqref{eq:symmetry} give
\begin{equation}
 K^\vee=(-1_P)^*D(K)\simeq K.
\label{eq:self-duality}
\end{equation}
This is the self-duality used below; the same pairing reappears in Section~7.

\begin{lemma}[The chosen twist is tensor-compatible]
\label[lemma]{lem:tensor-compatible-twist}
For every $\alpha\in\Omega$, the functor
\[
 \omega_\alpha(Q)=H^0(P,Q\otimes\mathcal L_\alpha)
\]
on the semisimple tensor category generated by the image of $K=\IC_X$ in
$\overline{\operatorname{Perv}}(P)$ is a fibre functor.  In particular the
Hodge decomposition of \cref{prop:hodge-vector} defines the Hodge
cocharacter on the same Tannakian category.
\end{lemma}

\begin{proof}
The decomposition
$Rj_*\QQ_S[2]\simeq K\oplus\delta_0$ and the definition of $\Omega$ give
\[
 H^i(P,K\otimes\mathcal L_\alpha)=0\qquad(i\ne0).
\]
Let $m_n:P^n\to P$ be addition.  Since a rank-one character local system
satisfies
$m_n^*\mathcal L_\alpha\simeq\mathcal L_\alpha^{\boxtimes n}$, the
projection formula and K\"unneth give
\begin{equation}
 H^\bullet(P,K^{*n}\otimes\mathcal L_\alpha)
 \simeq
 H^\bullet(P,K\otimes\mathcal L_\alpha)^{\otimes n},
\label{eq:tensor-twist-kunneth}
\end{equation}
so the left-hand side is concentrated in degree zero.

The complexes $K^{*n}$ are projective direct images of pure polarizable
Hodge modules, hence the decomposition theorem splits each of their perverse
cohomology sheaves as direct summands.  Consequently
${}^pH^0(K^{*n})$, and every simple direct summand of it, has twisted
cohomology concentrated in degree zero.  All other perverse-cohomological
pieces are negligible after passage to
$\overline{\operatorname{Perv}}(P)$, because convolution is $t$-exact in
the quotient.  Moreover, any negligible simple summand of
${}^pH^0(K^{*n})$ has Euler characteristic zero; the preceding concentration
then forces its twisted $H^0$ to vanish.  Hence $H^0(-\otimes\mathcal
L_\alpha)$ factors through the quotient on the tensor category under
consideration.  Since $K^\vee\simeq K$, every object in the semisimple tensor
category generated by $K$ is a direct summand of a finite direct sum of such
${}^pH^0(K^{*n})$.  Thus the same vanishing holds for every object of that
category.

Equation \eqref{eq:tensor-twist-kunneth}, by naturality and passage to direct
summands, supplies the tensor constraint for $\omega_\alpha$; exactness and
faithfulness follow from semisimplicity and the equality
$\dim H^0(P,Q\otimes\mathcal L_\alpha)=\chi(Q)>0$ for every nonzero simple
object in the quotient.  This agrees with the generic-twist fibre-functor
formalism of \cite[Section~3.1, Proposition~3.1]{KLM2026}.
\end{proof}


\section{The generic pair-incidence splitting}\label{sec:pair-splitting}

Return to the family over $\cB$.  Let $\cE\to\cB$ be the Weierstrass elliptic
scheme with zero section $O$, put $\eta=\mathcal O_{\cE}(O)$, and let
\[
 \pi:\cC\longrightarrow\cE,
 \qquad
 \kappa=\pi^*\eta
\]
be the relative shifted cyclic cover and its theta characteristic.  Let
$\cP\to\cB$ be the relative Prym scheme and $\cX\subset\cP$ the relative
centred theta section.  For any $Y\to\cB$, undecorated symbols
$E_Y,C_Y,P_Y,X_Y,\eta$ and $\kappa$ denote the corresponding base changes.

The relative theta divisor determined by $\kappa$ restricts to a section of
a relatively ample line bundle on the abelian scheme $\cP\to\cB$.  Its
restriction to every closed fibre is the nonzero Cartier divisor constructed
in \cref{prop:theta-resolution}.  Hence its zero scheme
$\cX\subset\cP$ is a symmetric relative effective Cartier divisor; the
defining section is a nonzerodivisor on every fibre of the smooth morphism
$\cP\to\cB$, so $\cX\to\cB$ is flat.

The fibrewise assertions proved above for
$b\in\cB(\CC)$ extend to all geometric fibres without an implicit
genericity step.  The locus of $b\in\cB$ for which $X_b$ is geometrically
integral is open for this proper flat family
\cite[IV$_3$, Th\'eor\`eme~12.2.4\textup{(viii)}]{EGA4}; it contains every
closed point by \cref{prop:theta-resolution}, hence equals $\cB$ because
$\cB$ is Jacobson.  Likewise, the nonsmooth locus of
$(\cX\setminus0)\to\cB$ is closed and would contain a closed point if
nonempty, while smoothness along the zero section would be an open condition
and would again contain a closed point if it occurred anywhere.  Thus every
geometric fibre is integral with singular locus exactly the origin, and the
relative smooth locus is
\[
 \cM=(\cX/\cB)_{\mathrm{sm}}=\cX\setminus0.
\]
The morphism $\cM\to\cB$ is smooth and surjective with geometrically
integral fibres.  Since $\cB$ is integral, $\cM$ is integral.

Let $\cN$ be the relative norm fibre in $\cC^{(3)}$.  Every Abel class on
$\cM$ has $h^0=1$, and the relative Abel map identifies
$\cN\times_{\cX}\cM$ with $\cM$.  The universal divisor on the relative
symmetric cube therefore gives a universal effective divisor of degree three
over $\cM$.

\subsection{A dominant normalized chart}

\begin{lemma}[Dominant normalized chart]\label[lemma]{lem:dominant-chart}
There is a nonempty open subscheme
\[
 T_{\QQ}\subset
 \operatorname{Spec}\QQ[c,c^{-1},m,m^{-1},e_1,e_2,e_3],
\]
hence an integral five-dimensional rational scheme, such that with
$\lambda=m^{-3}$, after base change to $\CC$ suitable dense opens of
\[
 T=T_{\QQ}\times_{\QQ}\CC
 \qquad\text{and}\qquad
 \cM
\]
are isomorphic.  On these opens the universal divisor is described by
\eqref{eq:chart-curve-divisor} below.  On $e_1\ne0$, put
\begin{equation}
 \delta=\frac{c}{e_1^3},\qquad
 \Lambda=e_1^3\lambda,\qquad
 E_2=\frac{e_2}{e_1^2},\qquad
 E_3=\frac{e_3}{e_1^3},\qquad
 Z=e_1z.
\label{eq:normalized-variables}
\end{equation}
The residual-octic family and its factor-ordering cover are defined over
$T_{\QQ}$ and are birationally pulled back from the normalized chart
$e_1=1$.
\end{lemma}

\begin{proof}
Let $D$ be the unique reduced degree-three divisor representing a point of a
dense open in $\cM$.  Its norm on $E$ is a unique nonvertical line
\[
 y=mx+n,
 \qquad m\ne0.
\]
Set
\[
 x'=mx+n-c,
 \qquad
 Y'=y-c,
 \qquad
 w'=w.
\]
Then the norm line is $Y'=x'$.  If the three selected lifts have
$w'$-coordinates $\alpha,\beta,\gamma$, let $e_i$ be their elementary
symmetric functions and write
\[
 p(W)=W^3-e_1W^2+e_2W-e_3.
\]
Newton's identities give
\[
 \begin{aligned}
 \alpha^3+\beta^3+\gamma^3
   &=e_1^3-3e_1e_2+3e_3=: \Sigma_1,\\
 (\alpha\beta)^3+(\alpha\gamma)^3+(\beta\gamma)^3
   &=e_2^3-3e_1e_2e_3+3e_3^2=: \Sigma_2.
 \end{aligned}
\]
Consequently
\[
 \prod_{\xi\in\{\alpha,\beta,\gamma\}}(x'-\xi^3)
 =(x')^3-\Sigma_1(x')^2+\Sigma_2x'-e_3^3.
\]
Since $\lambda=m^{-3}$, the curve and the selected divisor are reconstructed
in the primed coordinates by
\begin{equation}
 \begin{aligned}
 (Y'+c)^2
   &=G(x')\\
   &:=(x'+c)^2+\lambda\bigl((x')^3-\Sigma_1(x')^2
                  +\Sigma_2x'-e_3^3\bigr),\\
 D&:\quad x'=Y'=(w')^3,\qquad p(w')=0.
 \end{aligned}
\label{eq:chart-curve-divisor}
\end{equation}
Conversely these formulas recover the original short-Weierstrass model.
Indeed write
\[
 g_2=1-\lambda\Sigma_1,
 \qquad
 g_1=2c+\lambda\Sigma_2,
 \qquad
 g_0=c^2-\lambda e_3^3,
 \qquad
 n_0=-\frac{g_2}{3\lambda}.
\]
Since $m^3\lambda=1$, the substitution $x'=mx+n_0$ gives
\[
 (Y'+c)^2=x^3+Ax+B,
\]
where
\[
 A=m\left(g_1-\frac{g_2^2}{3\lambda}\right),
 \qquad
 B=G(n_0).
\]
The original norm line is then $y=mx+n_0+c$.  These constructions are
inverse on dense opens.  They are defined over $\QQ$, and therefore
construct $T_{\QQ}$ with
\begin{equation}
 \CC(T)=\CC(\cM).
\label{eq:chart-function-field}
\end{equation}

The equations are equivariant for
\[
 (x',Y',w',c)\longmapsto(r^3x',r^3Y',rw',r^3c),
\]
which in the original coordinates is the usual Weierstrass scaling
\[
 (x,y,w,c,A,B)
 \longmapsto
 (r^2x,r^3y,rw,r^3c,r^4A,r^6B).
\]
Thus
\[
 (c,\lambda,e_1,e_2,e_3,m)
 \longmapsto
 (r^3c,r^{-3}\lambda,re_1,r^2e_2,r^3e_3,rm).
\]
The basis of $H^0(C,3\kappa-D)$ used below has weights $3,3,4$, so its
projective coordinates transform by
\begin{equation}
 [u:v:z]\longmapsto[u:v:r^{-1}z].
\label{eq:section-scaling}
\end{equation}
Taking $r=e_1^{-1}$ gives \eqref{eq:normalized-variables} and
\[
 \QQ(c,\lambda,e_1,e_2,e_3)
 =\QQ(\delta,\Lambda,E_2,E_3)(e_1).
\]
If $q_0=m/e_1$, then
\begin{equation}
 \QQ(T_{\QQ})
 =\QQ(\delta,\Lambda,E_2,E_3)(q_0,e_1),
 \qquad q_0^3=\Lambda^{-1}.
\label{eq:cubic-chart-field}
\end{equation}
This gives the normalized description used below.
\end{proof}

\subsection{The norm conic}

All formulas in this subsection are first established over the rational
function fields indicated below.  Put $Y=y-c=w^3$.  Since
$\deg(3\kappa)=9>2g(C)-2$, Riemann--Roch gives
$h^0(C,3\kappa)=6$.  For the evaluation map
\[
 H^0(C,3\kappa)\longrightarrow H^0(D,(3\kappa)|_D),
\]
the kernel has the following basis:
\[
 s_0=x-Y,
 \qquad
 s_1=Y-e_1w^2+e_2w-e_3,
 \qquad
 s_2=(x-e_3)w-e_1Y+e_2w^2.
\]
For later use in the norm calculation, the basis argument is included.  On
$D$ we have $x=Y=w^3$ and
$p(w)=w^3-e_1w^2+e_2w-e_3=0$, so all three sections vanish on $D$.  Moreover
$\deg(3\kappa-D)=6$ and $K_C=2\kappa$, hence Riemann--Roch gives
\[
 h^0(C,3\kappa-D)-h^0(C,D-\kappa)=3.
\]
On the chart under consideration
$[\mathcal O_C(D)\otimes\kappa^{-1}]\in X_{\mathrm{reg}}=X\setminus\{0\}$,
so the degree-zero line bundle $\mathcal O_C(D-\kappa)$ is nontrivial and
has no section.  Thus $h^0(C,3\kappa-D)=3$.  Finally, in the expansion over
$k(E)$ with basis $1,w,w^2$, the coefficient of $x$ in the $w$-part first
forces the coefficient of $s_2$ in a constant linear relation to vanish;
the coefficient of $x$ in the $1$-part then forces that of $s_0$ to vanish,
and the coefficient of $s_1$ follows.  The three displayed sections are
therefore a basis.

For $s=us_0+vs_1+zs_2$, write
\[
 s=\alpha+\beta w+\gamma w^2,
\]
where
\[
 \begin{aligned}
 \alpha&=u(x-Y)+v(Y-e_3)-e_1zY,\\
 \beta&=e_2v+z(x-e_3),\\
 \gamma&=-e_1v+e_2z.
 \end{aligned}
\]
Since $w^3=Y$, the cyclic norm is
\begin{equation}
 \Nm(s)=\alpha^3+\beta^3Y+\gamma^3Y^2-3\alpha\beta\gamma Y.
\label{eq:cyclic-norm}
\end{equation}
It is divisible on $E$ by the fixed norm line $Y-x$.  The quotient is a
section of $\mathcal O_E(6O)$ and hence has a unique expression
\[
 Q_s=q_{00}+q_{01}x+q_{02}Y+q_{11}x^2+q_{12}xY+q_{22}Y^2.
\]
Let
\[
 M_s=
 \begin{pmatrix}
 q_{00}&q_{01}/2&q_{02}/2\\
 q_{01}/2&q_{11}&q_{12}/2\\
 q_{02}/2&q_{12}/2&q_{22}
 \end{pmatrix}.
\]

\begin{proposition}[Determinant factorization]
\label[proposition]{prop:determinant-factorization}
Put
\[
 K_0=\QQ(c,\lambda,e_1,e_2,e_3).
\]
In $K_0[u,v,z]$ one has
\begin{equation}
 \det(M_s)=(u-v)F_8,
\label{eq:determinant-factorization}
\end{equation}
where $F_8$ is homogeneous of degree eight and $u-v\nmid F_8$.
\end{proposition}

\begin{proof}
Each coefficient $q_{ij}$ is homogeneous of degree three in $(u,v,z)$, so
$\det(M_s)$ is homogeneous of degree nine.  The structural factor can be
seen without expanding the determinant.  Put
\[
 h=x-e_1w^2+e_2w-e_3.
\]
The basis above satisfies
\[
 s_0+s_1=h,\qquad s_2=wh.
\]
Hence on the hyperplane $u=v$ one has
\[
 s=h(u+zw).
\]
The section $h\in H^0(C,2\kappa)$ vanishes on the fixed divisor $D$.
Write $\operatorname{div}(h)=D+D'$.  Then
\[
 [\mathcal O_C(D')\otimes\kappa^{-1}]=-x,
 \qquad
 \Nm_\pi\mathcal O_C(D')=\eta^3.
\]
Thus $\pi_*D'$ is a line section of $E$, and the norm of $h$ is, up to a
nonzero scalar, the product of the fixed norm line $Y-x$ and this second
line.  Moreover
\[
 \Nm_\pi(u+zw)=u^3+z^3Y.
\]
After dividing by the fixed line, $Q_s$ is therefore a product of two lines
whenever $u=v$.  Its symmetric matrix has rank at most two, and
$\det(M_s)$ vanishes identically on $u=v$.  Thus $u-v$ divides the
determinant, and the quotient is homogeneous of degree eight.

It remains to show that this factor is simple.  Divisibility of $F_8$ by the
monic linear form $u-v$ is equivalent to the identity
$F_8(v,v,z)=0$ over $K_0$, so one specialization in characteristic zero at
which $F_8(v,v,z)$ is nonzero suffices.  Set
\[
 (c,\lambda,e_1,e_2,e_3)=(1,1,1,0,2),
 \qquad (u,v,z)=(1+t,1,1).
\]
The normalized curve is
\[
 Y^2+2Y=x^3-6x^2+14x-8.
\]
The six conic coefficients and the exact division by $Y-x$ are written in
Appendix~A.1.  Substitution into the determinant gives
\begin{equation}
 4\det M_{1+t,1,1}=2t(24t^3+56t^2+33t+6).
\label{eq:characteristic-zero-simple-factor}
\end{equation}
Hence $4F_8(1,1,1)=12\ne0$ in $\QQ$, so $u-v\nmid F_8$.
\end{proof}

\begin{lemma}[Geometric integrality of the residual octic]
\label[lemma]{lem:residual-octic-integral}
The projective curve
\[
 \mathscr R=V(F_8)\subset\PP^2_{K_0}
\]
is geometrically integral.
\end{lemma}

\begin{proof}
Reduce the normalized specialization
\[
 (c,\lambda,e_1,e_2,e_3)=(1,1,1,0,2)
\]
modulo $13$.  The corresponding short-Weierstrass curve is
\[
 y^2=x^3+2x+5,
\]
whose discriminant is nonzero; the branch cubic obtained from $y=1$ is also
separable.  Multiplying $F_8$ by the nonzero scalar $4$, set
\[
 H(u,v)=4F_8(u,v,1)
\]
at this specialization.  It has total degree eight and $v$-degree eight,
with constant leading coefficient $7$.  At $u=4$ its monic slice is
\[
 \begin{aligned}
 h(v)={}&v^8-3v^7-2v^6+3v^5+4v^4+4v^3\\
       &-3v^2-5v-2\in\mathbf F_{13}[v].
 \end{aligned}
\]
Appendix~A.3 records the Frobenius reductions
\[
 v^{13^8}\equiv v\pmod h,
 \qquad
 \gcd(h,v^{13^4}-v)=1,
\]
together with an explicit B\'ezout identity for the second relation.  The
standard factor-degree criterion then makes $h$ irreducible over
$\mathbf F_{13}$.
If $H$ factored in $\mathbf F_{13}[u,v]$, the constant leading coefficient in
$v$ would force the leading $v$-coefficients of the factors to be units.  A
positive-$v$-degree factor would retain its degree after $u=4$, while a
$v$-degree-zero factor would itself be a unit.  Both possibilities contradict
the irreducibility of $h$.  Hence $H$ is irreducible.

Let $F_{8,\mathrm{sp}}(u,v,z)$ denote the homogeneous specialization of
$F_8$.  Its dehomogenization at $z=1$ is $H/4$, and the nonzero $v^8$ term
shows that $z\nmid F_{8,\mathrm{sp}}$.  If
$F_{8,\mathrm{sp}}=AB$ were a nontrivial homogeneous factorization, then
neither $A(u,v,1)$ nor $B(u,v,1)$ could be a unit: a positive-degree
homogeneous polynomial whose dehomogenization is constant is a scalar
multiple of a power of $z$, which would force $z$ to divide
$F_{8,\mathrm{sp}}$.  Setting $z=1$ would therefore give a nontrivial
factorization of $H$, a contradiction.  Thus the projective special fibre is
irreducible.

The affine curve $H=0$ contains the rational point $(u,v)=(1,6)$, where
\[
 (\partial H/\partial u,\partial H/\partial v)=(2,0).
\]
Thus the point is smooth.  Since $\mathbf F_{13}$ is perfect, the irreducible
special fibre is geometrically reduced.  If it were not geometrically
irreducible, Frobenius would permute its geometric components transitively.
A rational point would then lie on every component and could not be a smooth
point.  The special fibre is therefore geometrically integral.

Spread the coefficients over a localization $R_0=\mathbf Z[1/N]$ with
$13\nmid N$, and choose an affine open
\[
 \mathscr U\subset\mathbb A^5_{R_0}
 =\operatorname{Spec}R_0[c,\lambda,e_1,e_2,e_3]
\]
containing the displayed characteristic-$13$ point on which the required
denominators and one coefficient of $F_8$ are units.  After multiplying by a
unit, $F_8$ defines a flat projective relative Cartier divisor of degree eight
in $\PP^2_{\mathscr U}$.  Geometric integrality is open in such a family
\cite[IV$_3$, Th\'eor\`eme~12.2.4\textup{(viii)}]{EGA4}.  Since the displayed
special fibre is geometrically integral, so is the generic fibre
$\mathscr R$.
\end{proof}

\begin{lemma}[Normalized and full function fields]
\label[lemma]{lem:normalized-function-field}
Work over $e_1\ne0$ and put
\[
 K_{\mathrm n}=\QQ(\delta,\Lambda,E_2,E_3).
\]
On the affine section chart $v\ne0$, put
\[
 U=u/v,\qquad z_0=z/v,\qquad \zeta=e_1z_0.
\]
Let $\mathscr R_{\mathrm n}$ be obtained by setting $e_1=v=1$ and using
coordinates $(\delta,\Lambda,E_2,E_3;U,\zeta)$.  Under the parameter map of
\cref{lem:dominant-chart}, let
\[
 \mathscr R_{T_{\QQ}}
 =\mathscr R\times_{\operatorname{Spec}K_0}
   \operatorname{Spec}\QQ(T_{\QQ}).
\]
Then
\begin{equation}
 k(\mathscr R)=k(\mathscr R_{\mathrm n})(e_1),
 \qquad
 k(\mathscr R_{T_{\QQ}})
 =k(\mathscr R_{\mathrm n})(e_1,q_0),
 \quad q_0^3=\Lambda^{-1}.
\label{eq:residual-function-fields}
\end{equation}
The curve $\mathscr R_{\mathrm n}$ is geometrically integral.  After base
change to $\CC$, the factor-ordering cover on the full chart is the pullback
of the normalized one through these field extensions.
\end{lemma}

\begin{proof}
The scaling in \cref{lem:dominant-chart}, together with
\eqref{eq:section-scaling}, identifies the locus $e_1v\ne0$ birationally with
$\GGm\times\mathscr R_{\mathrm n}$.  This gives the first equality in
\eqref{eq:residual-function-fields}.  The chart $v\ne0$ is dense because
$\mathscr R$ is a geometrically integral octic.  Geometric integrality of
$\mathscr R_{\mathrm n}$ follows by faithfully flat descent from
\cref{lem:residual-octic-integral}.  The second equality is
\eqref{eq:cubic-chart-field}.  Formation of the two line factors of a
rank-two conic commutes with field extension, which proves the last assertion.
\end{proof}

\begin{proposition}[Cyclotomic ordering cover]
\label[proposition]{prop:ordering-cover}
On the normalized chart $e_1=v=1$, rename
\[
 (\delta,\Lambda,E_2,E_3,U,\zeta)
 \quad\text{as}\quad
 (c,\lambda,e_2,e_3,u,z)
\]
and put
\[
 K_1=\QQ(c,\lambda,e_2,e_3,u).
\]
Let
\[
 \widetilde F(z)=F_8(u,1,z),
 \qquad
 \ell=\operatorname{lc}_z(\widetilde F),
 \qquad
 F=\ell^{-1}\widetilde F.
\]
Then $\deg_z\widetilde F=6$, so $\ell\ne0$ and $F$ is monic of degree six.  If
\[
 \Delta=-4\left(q_{00}q_{11}-(q_{01}/2)^2\right),
\]
there are polynomials $A_{\mathrm c},B_{\mathrm c}\in K_1[z]$ with
$\deg A_{\mathrm c}=3$ and $\deg B_{\mathrm c}\le2$ such that
\begin{equation}
 F=A_{\mathrm c}^2+3\Delta B_{\mathrm c}^2.
\label{eq:norm-square-identity}
\end{equation}
If
\[
 L=K_1[z]/(F)=k(\mathscr R_{\mathrm n}),
\]
then
\begin{equation}
 [\Delta]=[-3]\quad\text{in }L^\times/L^{\times2},
 \qquad
 L(\sqrt\Delta)=L(\sqrt{-3}).
\label{eq:cyclotomic-square-class}
\end{equation}
There is a dense normal open $\mathscr R_{\mathrm n}^\circ$ in the rank-two
locus for which the intrinsic factor-ordering cover is
\begin{equation}
 \operatorname{Ord}_{\mathrm{fac},\mathscr R_{\mathrm n}^\circ}
 \simeq
 \mathscr R_{\mathrm n}^\circ\times_{\QQ}
 \operatorname{Spec}\QQ(\sqrt{-3}).
\label{eq:cyclotomic-ordering-cover}
\end{equation}
The same statement holds after pullback to a dense normal open of
$\mathscr R_{T_{\QQ}}$.  In particular, the ordering cover splits over
$\CC$.
\end{proposition}

\begin{proof}
Appendix~\ref{app:exact-norm} gives $F$, $\Delta$, $A_{\mathrm c}$, and
$B_{\mathrm c}$ in the present coordinates.  Section~A.2 proves
\eqref{eq:norm-square-identity} coefficient by coefficient in the polynomial
ring $\mathbf Z[c,\lambda,e_2,e_3,u,\lambda^{-1}][z]$; the seven displayed
equalities there are the entire algebraic verification.
The normalized residual curve is geometrically integral by
\cref{lem:normalized-function-field}, so $F$ is irreducible and $L$ is a
field.  The leading terms in A.2 give $\deg A_{\mathrm c}=3$ and
$\deg B_{\mathrm c}=2$ generically.  Reducing
\eqref{eq:norm-square-identity} modulo $F$ gives
\[
 -3\Delta=(A_{\mathrm c}/B_{\mathrm c})^2\ne0
\]
in $L$, proving \eqref{eq:cyclotomic-square-class} and showing that the
generic conic has rank exactly two.

For a rank-two symmetric $3\times3$ matrix $M$, one has
$\operatorname{adj}(M)=c_0aa^t$ for a nonzero vector $a$.  Hence the ratio
of any two nonzero principal cofactors is a square, so the square class of a
nonzero principal minor is intrinsic.  To see directly that this is the
ordering character, restrict the conic to the coordinate line $Y=0$.
If the two factors have coefficient rows $(a_0,a_1)$ and $(b_0,b_1)$ in
$(1,x)$, the resulting binary quadratic has discriminant
$(a_0b_1-a_1b_0)^2$; changing the sign of its square root interchanges the
two factors.  Thus the intrinsic quadratic algebra is obtained by
adjoining a square root of $\Delta$.  After shrinking to a normal open on
which $\Delta$ and $B_{\mathrm c}$ are units,
write $\varrho^2=-3$.  The maps
\[
 s\longmapsto\frac{A_{\mathrm c}}{\varrho B_{\mathrm c}},
 \qquad
 \varrho\longmapsto\frac{A_{\mathrm c}}{B_{\mathrm c}s}
\]
give inverse isomorphisms between the intrinsic algebra
$\mathcal O[s]/(s^2-\Delta)$ and the constant algebra
$\mathcal O[\varrho]/(\varrho^2+3)$.  This proves
\eqref{eq:cyclotomic-ordering-cover}.  The field extensions in
\eqref{eq:residual-function-fields} preserve the construction, and over
$\CC$ the constant quadratic cover is split.
\end{proof}

\begin{remark}[Arithmetic fibres]
The quadratic field $\QQ(\sqrt{-3})=\QQ(\zeta_3)$ explains the term
``cyclotomic.''  After spreading the rank-two open over a suitable
$\mathbf Z[1/N]$ with $6\mid N$, the two line factors at an
$\mathbf F_q$-point are
individually defined over $\mathbf F_q$ precisely when $-3$ is a square in
$\mathbf F_q$, equivalently when $q\equiv1\pmod3$.  If $q\equiv2\pmod3$,
they are exchanged by Frobenius and become defined over $\mathbf F_{q^2}$.
\end{remark}


\subsection{The scheme-theoretic Abel inverse}

Over the chart $T$ of \cref{lem:dominant-chart}, let
$p:C_T\to T$ be the relative curve and let $D$ be its universal divisor.
After restricting $T$, the sheaf
$p_*\mathcal O_{C_T}(3\kappa-D)$ is locally free of rank three.  The
residual equation $F_8=0$ defines a relative octic
\[
 \mathscr R_T\subset
 \PP_T\bigl(p_*\mathcal O_{C_T}(3\kappa-D)\bigr).
\]
After a further shrinking of $T$, the section $F_8$ is nonzero on every
projective-space fibre.  Thus $\mathscr R_T\to T$ is a flat projective
relative Cartier divisor.  Its generic fibre is geometrically integral by
\cref{lem:residual-octic-integral,lem:normalized-function-field}.  Flatness
forces every irreducible component to dominate the integral base $T$, so the
integral generic fibre gives a single component.  A nilpotent section would
vanish generically and hence be $\mathcal O_T$-torsion, which flatness
excludes.  Thus $\mathscr R_T$ is integral.

Let $\mathscr R_T^{(2)}\subset\mathscr R_T$ be the open subset on which the
residual conic has rank exactly two.  The relative Hilbert scheme of line
components of the universal conic is then finite \'etale of degree two over
$\mathscr R_T^{(2)}$.  Choosing one component determines the other and hence
orders them; denote this cover by
\[
 \operatorname{Ord}_{\mathrm{fac}}
    \longrightarrow\mathscr R_T^{(2)}.
\]
Its points will be written $([s],\ell_1,\ell_2)$, with
$Q_s=\ell_1\ell_2$ up to a unit pulled back from the base.

\begin{definition}[Admissible residual-norm locus]
\label[definition]{def:admissible-locus}
Let $\mathscr G\subset\mathscr R_T^{(2)}$ be the open locus on which the
following conditions hold after pullback to
$\operatorname{Ord}_{\mathrm{fac}}$:
\begin{enumerate}[label=\textup{(\roman*)}]
\item the fixed divisor $D$ is reduced;
\item the fixed norm line and the two lines $\ell_i=0$ meet $E$ transversely,
      their degree-three line sections are pairwise disjoint, and the residual
      line sections avoid the branch divisor of $C\to E$;
\item above each residual point, $s$ vanishes on exactly one sheet of
      $C\to E$, and every such zero is simple;
\item the two resulting degree-three zero divisors have $h^0=1$, and their
      Abel classes lie in the relative smooth locus of $\cX$.
\end{enumerate}
These conditions define an open subset of
$\operatorname{Ord}_{\mathrm{fac}}$ invariant under interchanging
$\ell_1$ and $\ell_2$.  Its image is open because the finite \'etale cover is
open, and invariance makes it the full inverse image of that image.  This
defines $\mathscr G$.
\end{definition}

\begin{lemma}[Relative ordered-zero construction]
\label[lemma]{lem:scheme-abel-inverse}
Put
\[
 \operatorname{Ord}_{\mathrm{fac},\mathscr G}
 =\operatorname{Ord}_{\mathrm{fac}}
   \times_{\mathscr R_T^{(2)}}\mathscr G.
\]
Let
$\operatorname{Ord}_{\mathrm{div}}$ be the functor over $\mathscr G$ whose
$W$-points are ordered pairs of relative effective Cartier divisors
$(D_1,D_2)$ of degree three on $C_W/W$ such that
\[
 \operatorname{div}(s_W)=D_W+D_1+D_2
\]
and the unordered pair $\{\pi_*D_1,\pi_*D_2\}$ is the unordered pair of line
sections of $E_W$ defined by the two components of $Q_{s_W}$.  This functor is
represented by a finite \'etale double cover of $\mathscr G$.
The assignment $Q_s=\ell_1\ell_2\mapsto(D_1,D_2)$ is an isomorphism of finite
\'etale double covers
\[
 \operatorname{Ord}_{\mathrm{fac},\mathscr G}
 \xrightarrow{\sim}\operatorname{Ord}_{\mathrm{div}}.
\]
It commutes with arbitrary base change, including nonreduced base change.
Moreover, each $D_i$ is a relative effective Cartier divisor of degree three,
and $[\mathcal O(D_i)\otimes\kappa^{-1}]$ lies in $\cX_{\mathrm{sm}/\cB}$.
\end{lemma}

\begin{proof}
For the forward construction put
$T'=\operatorname{Ord}_{\mathrm{fac},\mathscr G}$ and write
$Q_s=\ell_1\ell_2$ in its universal ordering.  By definition of
$\mathscr G$, none of the tangency, branch, collision, multiple-zero, or Abel
speciality loci excluded in \cref{def:admissible-locus} occurs.  For the
remainder of the forward construction, every object with a subscript $T$ is
pulled back to $T'$; this base change is suppressed from the notation.
The line $\ell_i$ then
cuts out a relative finite \'etale divisor $B_i\subset E_T$ of degree three.
Moreover
\[
 U_i=C_T\times_{E_T}B_i\longrightarrow B_i
\]
is finite \'etale of degree three, because $B_i$ avoids the branch divisor.

The zero over $B_i$ is selected scheme-theoretically as follows.  \'Etale-locally on
$B_i$, split $U_i=B_i\sqcup B_i\sqcup B_i$ and trivialize the restriction
of $3\kappa$, using the induced trivialization of its norm line.  If
$a_1,a_2,a_3\in\mathcal O_{B_i}$ are the three values of $s$, compatibility
of the finite locally free norm with base change gives
\begin{equation}
 a_1a_2a_3=\Nm_\pi(s)|_{B_i}=0
\label{eq:split-norm-zero}
\end{equation}
as an equality in the structure sheaf, not merely at geometric points.
The collision exclusions imply that locally two values, say $a_2,a_3$, are
units.  Equation~\eqref{eq:split-norm-zero} then forces $a_1=0$, also in
nilpotent directions.
Consequently
\[
 Z_i=Z(s)\times_{E_T}B_i\subset U_i
\]
is locally exactly one open-and-closed sheet of $U_i/B_i$.  Thus
$Z_i\to B_i$ is finite \'etale of degree one and hence is an isomorphism.

Let $t$ be a local equation for $B_i$ in $E_T$.  Since $B_i$ occurs with
multiplicity one in the norm divisor and the other norm factors are
disjoint from it, $\Nm_\pi(s)=t u$ with $u$ a unit along $B_i$.  \'Etale-locally
on a neighbourhood of $B_i$ in the complement of the branch divisor, split
the cover $C_T\to E_T$.  On the selected sheet the other two values of $s$
remain units, so the same local product formula shows that the selected
value is $t$ times a unit.
The zero is therefore simple.  The image $D_i\subset C_T$ of $Z_i$ is
finite \'etale of degree three over $T'$; as a finite \'etale multisection of
the smooth relative curve $C_T/T'$, it is a relative effective Cartier
divisor.

The pairwise disjoint divisors $D,D_1,D_2$ lie in $Z(s)$ and have total
relative degree nine, equal to $\deg(3\kappa)$.  After dividing by their
Cartier equations, the residual section has degree zero and is nonvanishing
on every geometric fibre.  Nakayama's lemma makes it a trivialization, so
there is no vertical or nilpotent residual contribution and
\begin{equation}
 Z(s)=D+D_1+D_2
\label{eq:relative-zero-sum}
\end{equation}
as relative Cartier divisors.  Conversely, an ordered divisor datum
determines the unique projective section $[s]$ with zero divisor
$D+D_1+D_2$, and taking its norm recovers the ordered line sections $B_i$
and hence the ordered factors $\ell_i$.  Thus these assignments are inverse
already at the divisor level.

It remains to compare this divisor construction with the relative Abel
incidence.  Let
\[
 \alpha:C_T^{(3)}\longrightarrow\operatorname{Pic}^3(C_T/T)
\]
be the relative Abel map and let
$\mathcal F=\Nm^{-1}(\eta^3)\subset\operatorname{Pic}^3(C_T/T)$.  The norm fibre is
the cartesian base change
\[
 N_T=C_T^{(3)}\times_{\operatorname{Pic}^3(C_T/T)}\mathcal F.
\]
Let $W_3^{(1)}\subset\operatorname{Pic}^3(C_T/T)$ be the open on which $h^0=1$
and the pushforward of a Poincar\'e bundle is locally free of rank one and
commutes with base change.  The standard
projective-bundle description of the Abel map and cohomology and base change
then identify its source with the projectivization of that rank-one bundle.
Thus the map is an isomorphism and remains so after arbitrary base change.
Fibrewise, this is the $h^0=1$ case of Riemann--Kempf
\cite{Kempf1973}:
\[
 \alpha^{-1}(W_3^{(1)})\xrightarrow{\sim}W_3^{(1)}.
\]
Base changing this isomorphism by
$\mathcal F\to\operatorname{Pic}^3(C_T/T)$ and
translating by $\kappa^{-1}$ gives
\[
 N_T^{(1)}\xrightarrow{\sim}X_T^{(1)}.
\]
Here, by definition, $X_T^{(1)}$ is the open part of the theta section
represented by line bundles with $h^0=1$.
Restrict both sides to the inverse images of
$X_T^{(1)}\cap X_{T,\mathrm{sm}/T}$, as imposed above.  Thus the divisor
construction descends to the relative smooth Abel incidence.

Zero schemes, fibre products, finite locally free norms, and the cartesian
Abel construction all commute with arbitrary base change.  Hence both
composites preserve the universal zero schemes and are the identity as
morphisms.  The use of $C_T^{(3)}\simeq\operatorname{Hilb}^3(C_T/T)$ is only
representability; no normalization of the ordering cover or of pair fibres
is involved.
\end{proof}

Let $x_T:T\to\cM$ be the point represented by the universal divisor $D$,
so that $x_T=[\mathcal O(D)\otimes\kappa^{-1}]$.  Define
\[
 \cI_T^\dagger=
 \{y\in X_T^{(1)}\cap X_{T,\mathrm{sm}/T}:
       -x_T-y\in X_T^{(1)}\cap X_{T,\mathrm{sm}/T}\}.
\]
The two relative Abel inverses give universal divisors $D_1,D_2$ on
$C_T\times_T\cI_T^\dagger$.  Their Abel classes satisfy
\[
 [\mathcal O(D+D_1+D_2)\otimes\kappa^{-3}]=0
 \quad\text{in }\operatorname{Pic}^0(C_T/T).
\]
Thus $\mathcal O(D_1+D_2)$ and $\mathcal O(3\kappa-D)$ differ only by a line
bundle pulled back from the base.  Their divisor sections determine a unique
projective relative section
\[
 \vartheta:\cI_T^\dagger\longrightarrow
 \PP_T\bigl(p_*\mathcal O_{C_T}(3\kappa-D)\bigr).
\]
Taking norms shows that $\vartheta$ factors through $\mathscr R_T$.  Set
\[
 \cI_T^{\mathrm{adm}}=\vartheta^{-1}(\mathscr G).
\]

\begin{proposition}[Abel-incidence identification]
\label[proposition]{prop:abel-incidence}
There is an isomorphism over $\mathscr G$
\begin{equation}
 \Phi:\operatorname{Ord}_{\mathrm{fac},\mathscr G}
       \xrightarrow{\sim}\cI_T^{\mathrm{adm}},
 \qquad
 ([s],\ell_1,\ell_2)\longmapsto
 [\mathcal O(D_1)\otimes\kappa^{-1}].
\label{eq:abel-incidence-isomorphism}
\end{equation}
Interchanging $\ell_1$ and $\ell_2$ corresponds to the involution
$y\mapsto-x_T-y$.  The isomorphism commutes with arbitrary base change.
\end{proposition}

\begin{proof}
On $\operatorname{Ord}_{\mathrm{fac},\mathscr G}$, put
\[
 y_i=[\mathcal O(D_i)\otimes\kappa^{-1}]\in P_T.
\]
The divisor equality \eqref{eq:relative-zero-sum} gives
\begin{equation}
 x_T+y_1+y_2
  =[\mathcal O(D+D_1+D_2)\otimes\kappa^{-3}]=0.
\label{eq:abel-class-sum}
\end{equation}
Hence $y_2=-x_T-y_1$, and the ordered-zero construction of
\cref{lem:scheme-abel-inverse} defines the morphism $\Phi$.  The defining
conditions of $\mathscr G$ show that it lands in
$\cI_T^{\mathrm{adm}}$.

Conversely, let $y$ be a $W$-point of $\cI_T^{\mathrm{adm}}$, write
$p_W:C_T\times_T W\to W$ for the projection, put
$y_1=y$ and $y_2=-x_T-y$, and let $D_1,D_2$ be their relative Abel inverse
divisors.  Reading \eqref{eq:abel-class-sum} from right to left gives
\[
 \mathcal O(D_1+D_2)
 \simeq\mathcal O(3\kappa-D)\otimes p_W^*\mathcal N
\]
for a line bundle $\mathcal N$ on $W$.  The divisor section therefore gives
a unique projective section $[s]$ with
$\operatorname{div}(s)=D+D_1+D_2$; the factor pulled back from $W$ does not
affect this projective point.  Since $y_i\in P_T$,
\[
 \Nm_\pi\mathcal O(D_i)=\Nm_\pi(\kappa)=\eta^3.
\]
Thus the effective divisor $\pi_*D_i$ on $E_T$ is cut out by a unique
projective line $\ell_i$.  Multiplicativity of the finite locally free norm
gives, up to a
unit from $W$,
\[
 \Nm_\pi(s)=\ell_D\ell_1\ell_2,
 \qquad Q_s=\ell_1\ell_2,
\]
where $\ell_D$ is the fixed norm line of $D$.  Admissibility makes the two
residual factors distinct and orders them by the ordered pair $(D_1,D_2)$.
This constructs the inverse of \eqref{eq:abel-incidence-isomorphism}.

The relative Abel inverse on the $h^0=1$ locus, formation of zero divisors,
and finite locally free norms commute with arbitrary base change.  The two
composites preserve the universal Cartier divisors and hence are the identity
as morphisms.  Swapping $D_1,D_2$ gives $y_1\leftrightarrow y_2$, which is
the stated incidence involution.
\end{proof}

\begin{lemma}[Nonemptiness of the admissible locus]
\label[lemma]{lem:admissible-nonempty}
The open subscheme $\mathscr G$ of \cref{def:admissible-locus} is nonempty in
characteristic zero.  It is dense in $\mathscr R_T$, dominates $T$, and its
geometric generic fibre is a nonempty dense open of the geometric generic
residual curve.
\end{lemma}

\begin{proof}
For nonemptiness, use the following characteristic-$13$ point of the
normalized chart:
\[
 (c,\lambda,e_1,e_2,e_3)=(1,1,1,0,2),
 \qquad
 (u:v:z)=(3:5:1).
\]
In the shifted coordinates $Y=y-c$, the elliptic equation is
\[
 Y^2+2Y=x^3-6x^2+x+5
 \qquad(\bmod 13),
\]
equivalently, after the translation $x\mapsto x+2$,
\[
 y^2=x^3+2x+5.
\]
The cubic on the right has discriminant $8$, and the branch cubic cut out
by $Y=0$ has discriminant $4$.  The selected fixed divisor is defined by
\[
 p(w)=w^3-w^2-2,
\]
whose discriminant is $1$.  Thus both the curve and the branch divisor are
smooth and $D$ is reduced.

At the displayed section, the residual conic has coefficient vector
\[
 (q_{00},q_{01},q_{02},q_{11},q_{12},q_{22})=(5,1,7,7,0,2)
\]
and factors as
\[
 Q_s
 =5+x+7Y+7x^2+2Y^2
 =5(1+9x+Y)(1+12x+3Y).
\]
The fixed norm line is $Y=x$.  The discriminants of the three line sections
of $E$ are, respectively,
\[
 4,\qquad3,\qquad5,
\]
so all three intersections are transverse.  The two residual lines meet at
$(12,8)$; their intersections with $Y=x$ are $(9,9)$ and $(6,6)$.  None of
these three points lies on $E$, so the three degree-three line sections are
pairwise disjoint.  Each residual line avoids the branch divisor $Y=0$.

The canonical-model condition at this point follows from the same four forms
used in \cref{prop:theta-resolution}.  Since $\mathbf F_{13}$ contains the
cube roots of unity and $2,3$ are invertible, these forms remain regular.  The nine quadratic representatives other than
$X_0X_2-X_1^2$ split into the three deck-transformation eigenspaces
\[
 \langle1,x,Y,x^2\rangle,\qquad
 \langle w,xw,wY\rangle,\qquad
 \langle w^2,xw^2\rangle.
\]
At a point above $O$ their pole orders are, respectively,
$(0,2,3,4)$, $(1,3,4)$, and $(2,4)$.  Thus they are linearly independent,
and the unique canonical quadric is again
$X_0X_2-X_1^2=0$.  It has rank three and one ruling, so $|\kappa|$ is the
unique $g^1_3$.  In these coordinates a member of the pencil is cut out by
$a+bw$, and
\[
 \Nm_\pi(a+bw)=a^3+b^3Y.
\]
Thus its norm line has no $x$-term, whereas each of the three lines above has
a nonzero $x$-coefficient.  Their
degree-three inverse images are therefore not linearly equivalent to
$\kappa$ and hence lie in the $h^0=1$ Abel locus.  The local norm calculation
of \cref{lem:norm-fibre} is valid here for the same reason: $2$ and $3$ are
invertible.  At a residual point the cover is
\'etale and the residual line is transverse to $E$; hence the norm of $s$
has vanishing order one.  The three nonnegative vanishing orders of $s$ on
the three sheets sum to one, so exactly one sheet contains a simple zero.
This verifies all conditions in \cref{def:admissible-locus}.

For the passage to characteristic zero, enlarge the localization
$R_0=\mathbf Z[1/N]$ from \cref{lem:residual-octic-integral}, without
inverting $13$, so that the chart and all admissibility conditions spread
out.  Work on the affine chart $v\ne0$.  The coordinate ring $A_0$ of a
suitable affine open in the spread of $T_{\QQ}$ is a localization of
\[
 R_0[c,c^{-1},m,m^{-1},e_1,e_2,e_3]
\]
and hence is a unique factorization domain.  With
$U_0=u/v$ and $Z_0=z/v$, clear denominators in $F_8(U_0,1,Z_0)$ and divide
by its content.  The resulting primitive polynomial
$f\in A_0[U_0,Z_0]$ is irreducible over $\operatorname{Frac}(A_0)$ by
\cref{lem:residual-octic-integral,lem:normalized-function-field}; Gauss's
lemma makes it irreducible in $A_0[U_0,Z_0]$.  Thus
\[
 \mathscr Y=\operatorname{Spec}A_0[U_0,Z_0]/(f)
\]
is integral.  Since $R_0$ injects into this domain, its coordinate ring is
torsion-free as an $R_0$-module; over the localization $R_0$ of $\mathbf Z$
this is equivalent to flatness.  Hence $\mathscr Y$ is $R_0$-flat.

Admissibility defines an open subscheme of $\mathscr Y$ containing the
displayed point over $\mathbf F_{13}$.  It therefore contains a principal
open $D(g)$ with $g\ne0$.  Since $\mathscr Y$ is integral, $D(g)$ contains
the generic point of $\mathscr Y$; in particular, its characteristic-zero
generic fibre is nonempty.  The generic residual curve dominates the generic
point of $T_{\QQ}$, and this remains true on $D(g)$.  After base change to
$\CC$, this open lies in $\mathscr G$.  Hence $\mathscr G$ is nonempty, dense in $\mathscr R_T$, and
dominates $T$; its geometric generic fibre is a nonempty open of a
geometrically integral curve and is therefore dense.
\end{proof}

\begin{corollary}[The two admissible branches]
\label[corollary]{cor:two-admissible-branches}
After replacing $\mathscr G$ by a dense open subset, there is an isomorphism
over $\mathscr G$
\[
 \cI_T^{\mathrm{adm}}\simeq\mathscr G\sqcup\mathscr G.
\]
The incidence involution exchanges the two copies.  On a geometric generic
fibre over $T$, each copy is a dense open of a geometrically integral curve.
\end{corollary}

\begin{proof}
By \cref{prop:ordering-cover,lem:normalized-function-field}, the finite
\'etale factor-ordering cover has two rational sections over $\CC$ at the
generic point of $\mathscr G$.  On the principal-minor chart, after fixing
$\varrho\in\CC$ with $\varrho^2=-3$, they are given by
$s=\pm A_{\mathrm c}/(\varrho B_{\mathrm c})$ for the intrinsic coordinate
$s^2=\Delta$.
They extend to two disjoint sections after shrinking $\mathscr G$, so
$\operatorname{Ord}_{\mathrm{fac},\mathscr G}
 \simeq\mathscr G\sqcup\mathscr G$.
Apply \cref{prop:abel-incidence}.  The last assertion follows from
\cref{lem:residual-octic-integral,lem:admissible-nonempty}.
\end{proof}

\subsection{Boundary exhaustion and global labels}

Let $\xi$ be the generic point of the integral scheme $\cM$, and let
$\bar\xi$ be a geometric generic point.  For a geometric point $x$ of
$\cM$, define the pair fibre scheme-theoretically by
\[
 C_x=X\cap(-x-X)\subset P,
 \qquad C_x^{\mathrm{ne}}=C_x\setminus\{0,-x\}.
\]
Over $\bar\xi$, the unique effective divisors representing $x,y,-x-y$
satisfy
\begin{equation}
 D_x+D_y+D_{-x-y}=\operatorname{div}(s_y)
\label{eq:pair-zero-divisor}
\end{equation}
for a projective section $[s_y]\in\PP H^0(C,3\kappa-D_x)$.  This defines a
rational map
\[
 \Psi:C_x\dashrightarrow\PP H^0(C,3\kappa-D_x).
\]
\begin{lemma}[Boundary exhaustion]
\label[lemma]{lem:boundary-exhaustion}
 The rational map $\Psi$ has finite fibres wherever it is defined.  The
 determinant component $u=v$ gives no nonexceptional incidence-curve
 component; its generic divisor data give only the Abel classes $y=0,-x$.
 Every irreducible curve component of
$C_{\bar\xi}^{\mathrm{ne}}$ dominates the residual octic $\mathscr R$, meets
its admissible open, and is the closure of one of the two branches in
\cref{cor:two-admissible-branches}.  In particular,
$C_{\bar\xi}^{\mathrm{ne}}$ has exactly two irreducible components.
\end{lemma}

\begin{proof}
For a fixed section, the residual zero-cycle has only finitely many
degree-three subcycles, so $\Psi$ has finite fibres.  In the basis used for
\eqref{eq:determinant-factorization}, put
\[
 h=x-e_1w^2+e_2w-e_3.
\]
Since $Y=w^3$ on $C$, direct substitution gives
\begin{equation}
 s_0+s_1=h,\qquad s_2=wh.
\label{eq:structural-pencil}
\end{equation}
The zero divisor of $h\in H^0(C,2\kappa)$ is
$D_x+D_{-x}$.  Indeed, it contains $D_x$, and the residual degree-three
divisor has Abel class $-x$ and is unique on the generic $h^0=1$ locus.
Equation \eqref{eq:structural-pencil} identifies the line $u=v$ with
\[
 \PP\bigl(h\cdot H^0(C,\kappa)\bigr)
 =\PP H^0(C,3\kappa-D_x-D_{-x}).
\]
 Its members have divisor $D_x+D_{-x}+D_\kappa$ with
 $D_\kappa\in|\kappa|$.  For a general member of this pencil,
 $\pi_*D_{-x}$ and $\pi_*D_\kappa$ are reduced, disjoint line sections of the
 plane cubic $E$.  A
 degree-three subcycle using two points from the first section and one from
 the second cannot itself be a line section: a line through the first two
 points is the first line and cannot contain the third point.  The same
 argument applies with the roles reversed.  Hence the only degree-three
 subcycles whose pushforward is a line section are the two complete triples.
 Their Abel classes are $-x$ and $0$.

 The exceptional parameters of the structural pencil form a finite set.  If
 an incidence-curve component mapped dominantly to $u=v$, its generic point
 would contradict the preceding paragraph; if it mapped to one parameter, it
 would contradict the finiteness of the fibres of $\Psi$.  Thus $u=v$ gives no
 nonexceptional curve component.

 The same two-line argument shows that, over
 $\mathscr G$ the only Prym degree-three subcycles of the six residual zeros
are the two complete line sections.  By
\cref{prop:abel-incidence,cor:two-admissible-branches}, these give exactly
the two displayed admissible branches.

By the uniqueness of the $g^1_3$, the only points where an Abel inverse in
the definition of $\Psi$ is not unique are $y=0$ and $y=-x$.  Thus the
indeterminacy of $\Psi$ is contained in these two exceptional points.

Where $\Psi$ is defined, \eqref{eq:pair-zero-divisor} and multiplicativity
of the finite locally free norm give, up to a nonzero scalar,
\[
 \Nm_\pi(s_y)=\ell_{D_x}\ell_{D_y}\ell_{D_{-x-y}}.
\]
After division by the fixed line $\ell_{D_x}$, the residual conic is the
product $Q_{s_y}=\ell_{D_y}\ell_{D_{-x-y}}$.  Hence
$\operatorname{rank}M_{s_y}\le2$ and $\det M_{s_y}=0$.  The image of
$\Psi$ is therefore contained in
\[
 V(\det M_s)=V(u-v)\cup V(F_8).
\]
The preceding paragraph excludes $V(u-v)$ as the image of a
nonexceptional curve component.  Consequently every nonexceptional curve
component maps to the residual octic $V(F_8)$.  Its image is one-dimensional
because
$\Psi$ has finite fibres, and hence is dense because $\mathscr R$ is
geometrically integral.  It must meet the dense open $\mathscr G$.  An
irreducible curve cannot contain both disjoint admissible branches as open
subsets, so their closures are exactly the two nonexceptional components.
Here $\CC(T)=\CC(\cM)$ from \eqref{eq:chart-function-field} identifies the
two geometric generic fibres.
\end{proof}

\begin{proposition}[The geometric generic pair fibre]
\label[proposition]{prop:generic-pair-fibre}
The scheme $C_{\bar\xi}$ is reduced, and
$C_{\bar\xi}^{\mathrm{ne}}$ has exactly two geometrically integral
irreducible components $A_{\bar\xi}$ and
$B_{\bar\xi}$.  They are exchanged by $y\mapsto-x-y$, and both descend to
$k(\xi)$.
\end{proposition}

\begin{proof}
By
\cref{lem:residual-octic-integral,prop:ordering-cover,%
lem:admissible-nonempty,prop:abel-incidence,%
cor:two-admissible-branches,lem:boundary-exhaustion},
the nonexceptional geometric generic fibre has exactly two geometrically
integral components, exchanged by the incidence involution.  Since
$\Stab_P(X)=\{0\}$ and $x$ is generic, the integral Cartier divisors $X$ and
$-x-X$ are distinct.  Their local equations form a regular sequence in the
smooth threefold $P$, so $C_{\bar\xi}$ is a pure local complete intersection
curve.  It is therefore Cohen--Macaulay and has no embedded associated
points.  The two dense admissible branches show that it is reduced at every
generic point.  Hence $C_{\bar\xi}$ is reduced.

It remains to descend the two labels.  Let $\tau$
be the order-three deck transformation.  On $P$ one has
\[
 1+\tau+\tau^2=0,
\]
and $\tau(X)=X$.  Therefore
\[
 y_+(x)=\tau x\in C_x,
 \qquad s_+(x)=(x,\tau x)
\]
defines a rational section of the pair-incidence family; the associated
ordered pair is $(\tau x,\tau^2x)$.  This section meets the smooth locus of the projection of the pair-incidence
family to the $x$-coordinate on a nonempty open.  If it did not, then for
general $x$ the tangent hyperplanes to $X$ at $\tau x$ and $\tau^2x$ would
coincide, so their Gauss covectors would be equal.  By $\tau$-equivariance,
that common covector would be fixed projectively by $\tau$.  The general
Gauss image would then lie in the projective fixed locus of $\tau$, contrary
to the dominance in \eqref{eq:gauss-degree}.  Explicitly,
\[
 T_0^*P=
 \left\langle\frac{dx}{w(w^3+c)}\right\rangle
 \oplus
 \left\langle
   \frac{dx}{w^2(w^3+c)},
   \frac{x\,dx}{w^2(w^3+c)}
 \right\rangle,
\]
and the two summands have distinct $\tau$-eigenvalues.  The projective fixed
locus is a point union a line, a proper subset of $\PP^2$.

Over $k(\xi)$, the generic value of $y_+$ is a smooth rational point and
lies on exactly one geometric component.  Galois therefore fixes the
component containing $y_+$ and, since there are exactly two components,
also fixes the other.  Both components therefore descend to $k(\xi)$.
\end{proof}

\section{Spreading the generic splitting}\label{sec:spreading}

Define the relative pair-incidence scheme
\[
 \cI=\{(x,y)\in\cM\times_{\cB}\cX:-x-y\in\cX\},
\]
with the scheme structure induced by the two theta equations, and let
$q:\cI\to\cM$ be the first projection.  The equations $y=0$ and $y=-x$
define two sections of $q$.  Let $\cI^{\mathrm{ne}}$ denote their
complement.

\begin{theorem}[Uniform pair splitting]
\label{thm:uniform-pair-splitting}
There are a nonempty open subscheme $\cM^\circ\subset\cM$ and closed
subschemes $\cA_1,\cA_2\subset\cI|_{\cM^\circ}$ with the following
properties.
\begin{enumerate}[label=\textup{(\roman*)}]
\item Each $\cA_i\to\cM^\circ$ is flat, and every geometric fibre is a
      geometrically integral curve.
\item Every geometric fibre of $q|_{\cM^\circ}$ is reduced, and
\[
 \cI^{\mathrm{ne}}|_{\cM^\circ}
 = (\cA_1\cup\cA_2)\cap\cI^{\mathrm{ne}}
\]
      scheme-theoretically.
\item The two fibre curves are distinct at every geometric point of
      $\cM^\circ$, and the involution
      $(x,y)\mapsto(x,-x-y)$ exchanges $\cA_1$ and $\cA_2$.
\end{enumerate}
The image
\[
 \cB^\circ=\operatorname{im}(\cM^\circ\to\cB)
\]
is a nonempty Zariski-open subset.  For every $b\in\cB^\circ(\CC)$, the
fibre
\[
 U_b=\cM^\circ\times_{\cB}\{b\}
\]
is a nonempty dense connected open subset of $(X_b)_{\reg}$.  For every
$x\in U_b(\CC)$, the curves $(\cA_1)_x$ and $(\cA_2)_x$ are exactly the two
nonexceptional irreducible components of $C_x$.
\end{theorem}

\begin{proof}
Throughout the proof the schemes may be restricted to a smaller nonempty
open of $\cM$ without changing notation; the final open is denoted
$\cM^\circ$.
By \cref{prop:generic-pair-fibre}, the geometric generic fibre of
$q:\cI\to\cM$ is reduced and has exactly two geometrically integral
components, both defined over $k(\cM)$.  Let
$\cA_1,\cA_2\subset\cI$ be their closures.  The involution exchanges them
because it exchanges their generic fibres.

The morphism $q$ is projective and of finite presentation.  Apply generic
flatness to $q$ and to $\cA_i\to\cM$.  Geometric reducedness of the fibres of
$q$ and geometric integrality of the fibres of each $\cA_i$ then hold after
restricting to a nonempty open of $\cM$; this is the standard spreading of
these properties in a flat family \cite[\href{https://stacks.math.columbia.edu/tag/052A}{Tag 052A};
\href{https://stacks.math.columbia.edu/tag/0578}{Tag 0578};
\href{https://stacks.math.columbia.edu/tag/0559}{Tag 0559}]{stacks-project}.  This proves the
flatness and fibre assertions in (i), as well as the first assertion of (ii).

Since $q$ is projective and the $\cA_i$ are flat, their closed immersions
define two sections of the relative Hilbert scheme
$\operatorname{Hilb}_{\cI/\cM}$.  Its diagonal is closed.  The equalizer of
the two sections is therefore closed and omits the generic point.  Removing
it makes $(\cA_1)_x$ and $(\cA_2)_x$ distinct for every geometric $x$.

Give $\cA=\cA_1\cup\cA_2$ its scheme-theoretic union structure.  Let
$\mathcal J$ be the ideal of $\cA\cap\cI^{\mathrm{ne}}$ in
$\cI^{\mathrm{ne}}$.  The geometric generic fibre is reduced and consists
exactly of the two generic components, so $\mathcal J$ vanishes there.  The
image of $\operatorname{Supp}(\mathcal J)$ in the noetherian integral scheme
$\cM$ is constructible by Chevalley's theorem and does not contain its generic
point.  Its closure is therefore proper.  Remove this closure.  Then
$\mathcal J=0$, and
\[
 \cI^{\mathrm{ne}}=\cA\cap\cI^{\mathrm{ne}}
\]
over the remaining open.  The fibrewise reducedness, the distinctness and
geometric integrality of the two curve fibres, and this scheme-theoretic
equality prove (ii) and (iii).

Denote the remaining open by $\cM^\circ$.
The morphism $\cM^\circ\to\cB$ is open because it is the restriction of the
smooth morphism $\cM\to\cB$ \cite[\href{https://stacks.math.columbia.edu/tag/056G}{Tag 056G}]{stacks-project}.  Hence its image
$\cB^\circ$ is nonempty and open.  For $b\in\cB^\circ(\CC)$, the fibre
$U_b$ is a nonempty open of the integral surface $(X_b)_{\reg}$; it is
therefore dense, irreducible, and connected.  The fibrewise conclusions
(i)--(iii) give the final assertion.
\end{proof}

\section{Pair convolution detects \texorpdfstring{$\IC_X$}{IC(X)}}

Fix $b\in\cB^\circ(\CC)$, put $(P,X)=(P_b,X_b)$, and let $U=U_b$ as in
\cref{thm:uniform-pair-splitting}.
Let $\rho:S\to X$ be the minimal resolution and let $j:S\to P$ be its
composition with the inclusion.  Put
\[
 K=\IC_X,
 \qquad
 \widetilde K=Rj_*\QQ_S[2]=K\oplus\delta_0.
\]
Consider
\[
 \wt m:S\times S\longrightarrow P,
 \qquad (p,q)\longmapsto j(p)+j(q),
\]
and let $\sigma(p,q)=(q,p)$.  A superscript $-$ denotes the
anti-invariant direct summand under $\sigma$.

\begin{lemma}[Top cohomology on the split-pair locus]
\label[lemma]{lem:split-fibre-top}
After replacing $U$ by a dense open subset, put $U^+=-U$.  For $t=-x\in
U^+$, the fibre $F_t=\wt m^{-1}(t)$ has exactly four one-dimensional
irreducible components in its reduction: the strict transforms of
$(\cA_1)_x$ and $(\cA_2)_x$ and the two structural components
\[
 E_{\exc}\times\{s_t\},\qquad\{s_t\}\times E_{\exc}.
\]
Here $s_t$ is the unique point of $S$ above $t$.  The involution $\sigma$
exchanges the first two components and exchanges the last two.
Consequently
\[
 \dim\cH^{-2}(\widetilde K*\widetilde K)_{-,t}=2,
 \qquad
 \dim\cH^{-2}(K*K)_{-,t}=1.
\]
After a further shrinking of $U$, the restriction
$\cH^{-2}(K*K)_-|_{U^+}$ is the constant rank-one local system.
\end{lemma}

\begin{proof}
Fix $t\in U^+$, write $x=-t\in U$, and let
$E_{\exc}=\rho^{-1}(0)$.  By definition, $C_x$ parametrizes ordered pairs
whose sum is $-x=t$.  Hence the fibre
$F_t=\wt m^{-1}(t)$ has precisely four irreducible curve components: the
strict transforms of $(\cA_1)_x,(\cA_2)_x$ and the two structural components
\[
 C_{1,t}=E_{\exc}\times\{s_t\},
 \qquad
 C_{2,t}=\{s_t\}\times E_{\exc},
\]
where $s_t$ is the unique point above $t$.  Indeed, away from
$E_{\exc}$ the resolution is an isomorphism, while a pair with one entry on
$E_{\exc}$ must have the other entry over $t$.  Thus there are no further
top-dimensional components.  The involution $\sigma$ exchanges
$(\cA_1)_x$ with $(\cA_2)_x$ and $C_{1,t}$ with $C_{2,t}$.

Proper base change identifies
\[
 \cH^{-2}(\widetilde K*\widetilde K)_t=H^2(F_t,\QQ).
\]
The top cohomology of a proper complex curve is freely generated by the
fundamental classes of the irreducible components of its reduction.  Thus
$H^2(F_t,\QQ)_-$ has the explicit basis
\[
 [\,(\cA_1)_x\,]-[\,(\cA_2)_x\,],
 \qquad [C_{1,t}]-[C_{2,t}].
\]
Here and below the same notation is used for a curve and its strict transform
in $F_t$.  Holomorphic transposition preserves complex orientation, so
\begin{equation}
 \dim\cH^{-2}(\widetilde K*\widetilde K)_{-,t}=2.
\label{eq:total-alt-rank}
\end{equation}

Expand the decomposition $\widetilde K=K\oplus\delta_0$:
\[
 \widetilde K*\widetilde K=(K*K)\oplus(K*\delta_0)\oplus(\delta_0*K)
      \oplus(\delta_0*\delta_0).
\]
At $t\ne0$, the last term has zero stalk.  The two middle terms are both
canonically $K$, and transposition exchanges them.  Since $t$ is a smooth
point of the surface $X$, their anti-invariant stalk in degree $-2$ is the
line generated by $[C_{1,t}]-[C_{2,t}]$.  Taking dimensions in this direct
sum decomposition gives
\begin{equation}
 \dim\cH^{-2}(K*K)_{-,t}=1.
\label{eq:principal-alt-rank}
\end{equation}
There is no Koszul sign: $K$ has geometric shift $[2]$, and
$(-1)^{2\cdot2}=1$.

Shrink $U^+$ once more and choose a topological stratification.  The strict
transforms of $\cA_1$ and $\cA_2$ are then flat over $U^+$; proper
stratified local triviality makes all four relative fundamental classes
locally constant.  They are globally labelled as well.  The first two come
from $\cA_1$ and $\cA_2$ in \cref{thm:uniform-pair-splitting}, while the
other two are
\[
 C_1=E_{\exc}\times s,\qquad C_2=s\times E_{\exc},
\]
where $s:U^+\to S$ is the unique lift of the inclusion $U^+\subset X_{\reg}$.
Thus monodromy cannot permute these four component classes.  Since complex
orientation fixes the sign of each fundamental class,
\[
 [\cA_1]-[\cA_2],\qquad [C_1]-[C_2]
\]
form a global frame of the anti-invariant rank-two local system
$\cH^{-2}(\widetilde K*\widetilde K)_-|_{U^+}$.

On $U^+$ the sheaf $K$ is $\QQ_{U^+}[2]$.  Hence the anti-invariant part of
$(K*\delta_0)\oplus(\delta_0*K)$ is the constant line generated by
$[C_1]-[C_2]$.  The direct-sum decomposition above then identifies
$\cH^{-2}(K*K)_-|_{U^+}$ with the complementary constant line.  This proves
the last assertion.
\end{proof}

\begin{lemma}[Vanishing away from the theta divisor]
\label[lemma]{lem:off-theta-vanishing}
There is a dense open subset $W\subset P\setminus X$ such that, for every
$t\in W$, the scheme $X\cap(t-X)$ is a smooth connected projective curve.
The involution $p\mapsto t-p$ acts as the identity on its top cohomology.  In
particular,
\[
 \cH^{-2}(K*K)_-|_W=0.
\]
\end{lemma}

\begin{proof}
The sum map $m|_{X\times X}:X\times X\to P$ is surjective because two
translates of the ample divisor $X$ always meet.  Moreover,
\[
 m\bigl((\{0\}\times X)\cup(X\times\{0\})\bigr)=X.
\]
Hence, for $t\notin X$, the fibre lies entirely in
$X_{\reg}\times X_{\reg}$.  Generic smoothness for the restricted sum map
therefore gives smoothness of the general fibre.  Since $t\notin X$ and
$\Stab_P(X)=\{0\}$, the translate $t-X$ does not contain $X$; its restriction
to the integral Cohen--Macaulay surface $X$ is therefore an effective ample
Cartier divisor.  Hartshorne's connectedness theorem, applied on $X$, shows
that this divisor, namely $X\cap(t-X)$, is connected
\cite[Corollary~III.7.9]{Hartshorne1977}.
The involution $p\mapsto t-p$ preserves its complex orientation and hence
acts by $+1$ on its one-dimensional top cohomology.  Therefore
\begin{equation}
 \cH^{-2}(K*K)_-=0
 \quad\text{on a dense open subset of }P.
\label{eq:generic-alt-vanishing}
\end{equation}
\end{proof}

\begin{lemma}[Extraction of the divisor-supported summand]
\label[lemma]{lem:extract-ic}
In the decomposition-theorem splitting of $(K*K)_-$, the constant rank-one
local system of \cref{lem:split-fibre-top} is the restriction of an unshifted
$\IC_X$ summand.  It cannot arise from a shifted simple summand with support
$P$, or from a summand with any other support.
\end{lemma}

\begin{proof}
The complex $(K*K)_-$ is a direct summand of the projective direct image of
a pure polarizable Hodge module.  By the decomposition theorem and
semisimplicity \cite{BBD1982,Saito1988}, it is a direct sum of shifted simple
intersection complexes.  A summand with proper support containing the
generic point of $U^+$ must have support $X$.  Such a simple summand has the
form $\IC_X(\mathcal V)[r]$, where $\mathcal V$ is an irreducible local
system on a dense smooth open $X_{\mathcal V}\subset X_{\reg}$.  Its
restriction to $U^+\cap X_{\mathcal V}$ is $\mathcal V[2+r]$, so contribution
to ordinary degree $-2$ forces $r=0$.

A full-support summand cannot account for the line in
\eqref{eq:principal-alt-rank}.  Choose a general point $a\in U^+$ that lies
on no singular stratum of any full-support simple summand except the divisor
stratum $X_{\reg}$.  A small analytic neighbourhood of $a$ is stratified
homeomorphic to a product of a ball in $X_{\reg}$ with a normal disk
$\Delta$ to $X$.  Hence a full-support simple perverse summand is locally
\[
 \mathcal L[2]\boxtimes j_{\Delta,!*}\mathcal W[1],
 \qquad j_\Delta:\Delta^*\hookrightarrow\Delta,
\]
for local systems $\mathcal L$ and $\mathcal W$.

Write $T$ for the local monodromy of $\mathcal W$.  On a disk the middle
extension has the explicit stalk
\[
 \cH^k\!\left(i_0^*j_{\Delta,!*}\mathcal W[1]\right)
 =\begin{cases}
   \ker(T-1),&k=-1,\\
   0,&k\ne-1,
  \end{cases}
 \qquad i_0:\{0\}\hookrightarrow\Delta.
\]
Indeed $Rj_{\Delta,*}\mathcal W[1]$ has stalk cohomology
$\ker(T-1)$ in degree $-1$ and $\operatorname{coker}(T-1)$ in degree $0$;
the middle extension removes exactly the quotient supported at the origin.
It follows that an unshifted full-support perverse summand contributes along
$X$ only in ordinary degree $-3$.  To contribute to ordinary degree $-2$
it must occur with shift $[-1]$.  But on the dense smooth open of $P$ the
same shifted summand is a nonzero local system $\mathcal E[3][-1]$, hence
also contributes to ordinary degree $-2$ there, contradicting
\eqref{eq:generic-alt-vanishing}.  Thus no full-support summand contributes
to the rank-one line on $U^+$.

The two relevant shifts are summarized by
\[
\begin{array}{c|c|c}
\text{full-support summand}&\text{general point of }P&\text{general point of }X\\ \hline
\IC_P(\mathcal E)&-3&-3\\
\IC_P(\mathcal E)[-1]&-2&-2.
\end{array}
\]
Supports of dimension at most one do not meet a generic point of $U^+$, and
a different irreducible divisor cannot contain the nonempty open $U^+\subset
X$.

The local system $\mathcal V$ has rank one because its contribution to
\eqref{eq:principal-alt-rank} has rank one.  Replace $U^+$ by
$U^+\cap X_{\mathcal V}$.  Its restriction there is the constant local
system found in \cref{lem:split-fibre-top}.  The complement
$X_{\mathcal V}\setminus U^+$ is a proper complex algebraic subset and
therefore has real codimension at least two.  Every loop in the smooth
variety $X_{\mathcal V}$ can be perturbed off this complement, so the
inclusion induces a surjection
\[
 \pi_1(U^+)\twoheadrightarrow\pi_1(X_{\mathcal V}).
\]
The monodromy representation of $\mathcal V$ is trivial on the source and
hence on $\pi_1(X_{\mathcal V})$.  Thus $\mathcal V$ is trivial and the
summand is $\IC_X$ itself.
\end{proof}

\begin{proposition}[Alternating support detector]
\label[proposition]{prop:alt-support}
In $\operatorname{Perv}(P,\QQ)$, $K=\IC_X$ is a direct summand of
${}^pH^0((K*K)_-)$.  After tensoring with $\CC$, its image in
$\overline{\operatorname{Perv}}(P)$ is nonzero and is a direct summand of
$\Alt^2(K)$.
\end{proposition}

\begin{proof}
The first assertion is \cref{lem:extract-ic}.  Apply the quotient functor to
$K\otimes_{\QQ}\CC$.  It is perverse $t$-exact, convolution is $t$-exact in
$\overline D(P)$, and all nonzero perverse-cohomological degrees of $K*K$ map
to zero.  Hence the
image of ${}^pH^0((K*K)_-)$ is the anti-invariant summand of the tensor square
of the image of $K$.  Because the support dimension is even, this is
$\Alt^2(K)$ under the fibre functor.  The exact quotient functor preserves
the direct summand supplied by \cref{lem:extract-ic}; that summand is nonzero
because $\chi(K)=7$.
\end{proof}

\section{Identification of the full Tannaka group}

Let
\[
 G=G(K),\qquad V=\omega(K).
\]
By \eqref{eq:chi-seven}, $\dim V=7$.  The self-duality
\eqref{eq:self-duality} gives a perfect evaluation pairing
$K*K\to\delta_0$.  The chosen symmetry
is centred at the origin, so no translated skyscraper occurs, and
$\dim X=2$ makes the pairing symmetric.  Equivalently, the nondegenerate
self-duality on the odd-dimensional space $V$ cannot be alternating, so its
bilinear form is necessarily symmetric.  Therefore the full group embeds
faithfully as
\begin{equation}
 G\hookrightarrow \mathrm{O}(V).
\label{eq:orthogonal}
\end{equation}

The simple perverse sheaf $K$ is nondivisible: a translation fixing it must
fix its support, and \eqref{eq:stabilizer} then makes the translation zero.
By \cite[Proposition~3.3]{JKLM2025},
\begin{equation}
 V|_{G^0}\quad\text{is irreducible}.
\label{eq:connected-irreducible}
\end{equation}
This result does not require the ambient abelian variety to be simple.

Choose $\alpha\in\Omega$ as in \cref{prop:hodge-vector}.  By
\cref{lem:tensor-compatible-twist}, $H^0(-\otimes\mathcal L_\alpha)$ is a
tensor fibre functor on the category generated by $K$.  Since any two fibre
functors over $\CC$ are noncanonically isomorphic, it may be identified
with $\omega$.  Let $G_P=G$, and let
$G_H$ be the Tannaka group of the pure Hodge module $\IC_X^H$, acting
faithfully on the same complex vector space $V$.
By \cite[Corollary~3.2]{KLM2026}, the Hodge decomposition is defined by a
cocharacter $h:\GGm^2\to G_H$.  The comparison theorem
\cite[Corollary~2.5]{KLM2026} gives
\[
 G_P\mathrel{\triangleleft}G_H,
 \qquad
 [G_P^0,G_P^0]=[G_H^0,G_H^0]=:G^*.
\]
Since $G_P^0\subset G_H^0$, equation
\eqref{eq:connected-irreducible} implies that $V|_{G_H^0}$ is irreducible.
Thus the connected centre $Z_H^0$ acts by scalar matrices.  As $G_H^0$ is
reductive, multiplication $Z_H^0\times G^*\to G_H^0$ is surjective with
finite kernel.

Put $\mu(z)=h(z,z^{-1})$.  We need the antidiagonal cocharacter in the
common derived group, not merely in the Hodge group; the determinant gives a
short way to see this.  Since $\GGm$ is connected,
$\mu(\GGm)\subset G_H^0$.  By \cref{prop:hodge-vector},
\[
 \det\mu(z)=z^{(-2)\cdot2+0\cdot3+2\cdot2}=1.
\]
The determinant is trivial on the semisimple group $G^*$ and therefore
descends to a character
\[
 \overline{\det}:G_H^0/G^*\longrightarrow\GGm.
\]
The quotient $G_H^0/G^*$ is the image of $Z_H^0$.  On $Z_H^0$ the
determinant is $c\mapsto c^7$; since the faithful scalar action embeds
$Z_H^0$ into $\GGm$, the induced character $\overline{\det}$ has finite
kernel.  Hence the image of the connected group $\GGm$ under $\mu$ in
$G_H^0/G^*$ is both connected and finite, and is therefore trivial.
Consequently
\[
 \mu(\GGm)\subset G^*,
\]
the common Hodge and perverse derived group.  Its weights on $V$ are
\begin{equation}
 -2,0,2\quad\text{with multiplicities }2,3,2.
\label{eq:hodge-weights}
\end{equation}

Applying the fibre functor to \cref{prop:alt-support} gives
\begin{equation}
 \operatorname{Hom}_G(V,\Lambda^2V)\ne0.
\label{eq:hom-exterior}
\end{equation}

\begin{proposition}[Connected Tannaka group]\label[proposition]{prop:connected-g2}
There is an isomorphism of pairs
\[
 (G^0,V|_{G^0})\simeq(G_2,V_7).
\]
\end{proposition}

\begin{proof}
By \eqref{eq:orthogonal} and \eqref{eq:connected-irreducible}, the connected
centre of $G^0$ acts by orthogonal scalars and is therefore trivial.  Thus
$G^0$ is semisimple.  An irreducible representation of a product of
semisimple groups is an external tensor product.  Since $7$ is prime and the
action is faithful, $G^0$ has one simple factor.

For completeness, the seven-dimensional possibilities can be read directly
from the Weyl dimension formula (compare the standard tables in
\cite{Dynkin1957,FultonHarris1991}).  If a dominant highest weight contains
a fundamental weight $\omega_i$, its Weyl dimension is at least that of
$V_{\omega_i}$, and it increases strictly after adding a nonzero dominant
weight.  The only cases requiring inspection are therefore the small-rank
types.  They give
\[
\begin{array}{c|c|c}
\text{type}&\text{highest weight}&\text{dimension}\\ \hline
A_1&6\omega_1&7\\
A_6&\omega_1,\ \omega_6&7\\
B_3&\omega_1&7\\
G_2&\omega_1&7.
\end{array}
\]
Here is the short exclusion of the remaining types.  For $A_r$ the
fundamental dimensions are $\binom{r+1}{i}$.  The only ranks below $A_6$
that could conceivably contribute dimension seven are $A_2,A_3,A_4,A_5$.
The Weyl formula gives, respectively, the dimensions below the next jump as
$3,3,6,6$ (then $8$), $4,4,6$ (then $10$), $5,5$ (then $10$), and $6,6$
(then $15$).  Thus none has a seven-dimensional simple module.  For
$r\ge7$ even the standard $A_r$-module has dimension at least $8$.
For $C_2$ the fundamental dimensions are $4,5$ and the next irreducible
dimension is $10$; for $C_3$ they are $6,14,14$ and
$\dim V_{2\omega_1}=21$; for $C_r$ with $r\ge4$ the standard module
already has dimension at least $8$.  In type $B_r$, $r\ge4$, the vector
module has dimension at least $9$, and in type $D_r$, $r\ge4$, it has
dimension at least $8$; the spin modules do not produce dimension $7$.
Finally, the smallest nontrivial irreducible dimensions of
$F_4,E_6,E_7,E_8$ are $26,27,56,248$, whereas $G_2$ has its
seven-dimensional fundamental module.  This leaves exactly the four cases
in the table.

The orthogonal self-duality of $V$ excludes type $A_6$.  In type $A_1$, the
central element $-I\in\mathrm{SL}_2$ acts trivially on $\Sym^6(\CC^2)$,
so the faithful image is $\mathrm{PGL}_2$; in type $B_3$, the
seven-dimensional vector representation has faithful image $\mathrm{SO}_7$.
The remaining possibilities are therefore $\mathrm{PGL}_2$ on
$\Sym^6(\CC^2)$, $\mathrm{SO}_7$ on its standard module, and $G_2$ on
$V_7$.  For $\mathrm{SO}_7$, the exterior
square of the standard module is the irreducible adjoint module, so
$\operatorname{Hom}_{\mathrm{SO}_7}(V,\Lambda^2V)=0$, contrary to
\eqref{eq:hom-exterior}.  For a nontrivial cocharacter of $\mathrm{PGL}_2$, the weights on
$\Sym^6(\CC^2)$ are seven distinct numbers.  With the primitive
cocharacter of the adjoint torus (so that the usual $\mathrm{SL}_2$ weights
are divided by two), they are
\[
 -3,-2,-1,0,1,2,3,
\]
and every other nonzero cocharacter multiplies them by a nonzero integer.
The trivial cocharacter has one weight of multiplicity seven.  Neither
possibility agrees with \eqref{eq:hodge-weights}.  Hence $G^0=G_2$.
\end{proof}

\begin{corollary}[The full Tannaka group]\label[corollary]{cor:full-g2}
There is an isomorphism of pairs
\[
 (G,V)\simeq(G_2,V_7).
\]
In particular, $G$ is connected; this identifies the full Tannaka group, not
only its identity component.
\end{corollary}

\begin{proof}
By \cref{prop:connected-g2}, $G^0\simeq G_2$ on $V_7$.  Since
$G^0\mathrel{\triangleleft}G\subset \mathrm{O}(V)$,
\[
 G\subset N_{\mathrm{O}(V)}(G_2)=G_2\times\{\pm I\}.
\]
Indeed, conjugation gives a homomorphism from the normalizer to
$\operatorname{Aut}(G_2)$.  Since $\operatorname{Out}(G_2)=1$, multiplication
by a suitable element of $G_2$ puts every normalizing element in the
centralizer.  Schur's lemma makes that centralizer scalar, and its intersection
with $\mathrm{O}(V)$ is $\{\pm I\}$.  If $G\ne G^0$, choose
$g\in G\cap(-G_2)$ and write $g=(-I)h$ with $h\in G_2=G^0$.  Then
$-I=gh^{-1}\in G$.  But $-I$ acts by $-1$ on $V$ and by $+1$ on
$\Lambda^2V$, contradicting \eqref{eq:hom-exterior}.  Hence $G=G^0=G_2$.
\end{proof}

\begin{proof}[Proof of \cref{thm:main}]
Let $\cB^\circ$ be the nonempty open subset obtained in
\cref{thm:uniform-pair-splitting}.  The relative divisor $\cX\subset\cP$ is
symmetric, relatively ample, effective Cartier, flat, and has geometrically
integral fibres by the construction at the beginning of
\cref{sec:pair-splitting}.

Fix $b\in\cB^\circ(\CC)$.  The polarization type and the geometric
properties of $(P_b,X_b)$ follow from \cref{prop:theta-resolution} and the
construction of $L$.  The locus $\Omega_b$ in \cref{prop:hodge-vector} is
nonempty; for any $\alpha_b\in\Omega_b$, that proposition gives Hodge
multiplicities $(2,3,2)$.  The two pair components furnished by
\cref{thm:uniform-pair-splitting} imply, through \cref{prop:alt-support},
that $\IC_{X_b}$ is a direct summand of $\Alt^2(\IC_{X_b})$.  Finally,
\cref{cor:full-g2} identifies the full Tannaka group and tautological module
with $(G_2,V_7)$.  Since $b$ was arbitrary in $\cB^\circ(\CC)$, the theorem
follows.
\end{proof}

\appendix

\section{Exact norm identities and arithmetic specializations}\label{app:exact-norm}

This appendix isolates the finite algebra used in
\cref{sec:pair-splitting}.  The full residual octic is omitted because it
plays no independent conceptual role, but every algebraic assertion used in
the proof is displayed here.  Each step is either an identity in a specified
polynomial ring or a finite calculation in $\mathbf F_{13}$ using the
reduction relation written out below.  Consequently every algebraic input to
\cref{sec:pair-splitting} can be checked from the displayed formulas; no
unrecorded computation is used.
The four items recorded are the normalized conic coefficients, the
square-class identity, geometric integrality of the residual octic, and an
admissible specialization.

\subsection*{A.1. Normalized conic coefficients}
On the normalized chart $e_1=v=1$ the shifted curve is
\begin{equation}
 Y^2+2cY
 =\lambda x^3+\bigl(1-\lambda(1-3e_2+3e_3)\bigr)x^2
 +\bigl(2c+\lambda(e_2^3-3e_2e_3+3e_3^2)\bigr)x
 -\lambda e_3^3.
\label{eq:appendix-shifted-curve}
\end{equation}
For the section $s=us_0+s_1+zs_2$ one has
\[
 \alpha=u(x-Y)+(Y-e_3)-zY,\qquad
 \beta=e_2+z(x-e_3),\qquad
 \gamma=-1+e_2z.
\]
Substitute these expressions into
$\operatorname{Nm}(s)=\alpha^3+\beta^3Y+\gamma^3Y^2-3\alpha\beta\gamma Y$,
reduce $Y^2$ using \eqref{eq:appendix-shifted-curve}, and divide by the fixed
line $Y-x$.  The quotient
$Q_s=q_{00}+q_{01}x+q_{02}Y+q_{11}x^2+q_{12}xY+q_{22}Y^2$ has
\begin{align*}
q_{00}&=\frac{2c+\lambda(e_2^3-3e_2e_3-3e_3^2u+3e_3^2)}{\lambda},\\
q_{01}&=\frac{1+\lambda(3e_2+3e_3u^2-3e_3-1)}{\lambda},\\
q_{11}&=-(u-1)(u^2+u+1),
\end{align*}
and
\begin{align*}
\lambda q_{02}={}&2cz^3+e_2^3\lambda z^3+3e_2^2\lambda uz-3e_2^2\lambda z
 -3e_2e_3\lambda uz^2-3e_2e_3\lambda z^3+3e_2e_3\lambda z^2\\
&-3e_2\lambda u+3e_2\lambda-3e_3\lambda u^2-3e_3\lambda uz
 +6e_3\lambda u+3e_3\lambda z-3e_3\lambda-\lambda+1,\\
\lambda q_{12}={}&3e_2\lambda uz^2+3e_2\lambda z^3-3e_2\lambda z^2
 +2\lambda u^3+3\lambda u^2z-3\lambda u^2-3\lambda uz\\
&-\lambda z^3+\lambda+z^3,\\
\lambda q_{22}={}&-\lambda u^3-3\lambda u^2z+3\lambda u^2-3\lambda uz^2
 +6\lambda uz-3\lambda u\\
&-\lambda z^3+3\lambda z^2-3\lambda z+\lambda+z^3.
\end{align*}
These six coefficients determine $M_s$ and hence $F_8$.  In particular, for
\[
 M_s=\begin{pmatrix}
 q_{00}&q_{01}/2&q_{02}/2\\
 q_{01}/2&q_{11}&q_{12}/2\\
 q_{02}/2&q_{12}/2&q_{22}
 \end{pmatrix}
\]
the only determinant expansion used below is
\begin{equation}
 4\det M_s=
 4q_{00}q_{11}q_{22}-q_{00}q_{12}^2-q_{11}q_{02}^2
 -q_{22}q_{01}^2+q_{01}q_{02}q_{12}.
\label{eq:appendix-determinant-expansion}
\end{equation}

For the simple-factor assertion in
\cref{prop:determinant-factorization}, specialize in characteristic zero to
\[
 (c,\lambda,e_1,e_2,e_3)=(1,1,1,0,2),
 \qquad (u,v,z)=(1+t,1,1).
\]
Then \eqref{eq:appendix-shifted-curve} becomes
\[
 Y^2+2Y=x^3-6x^2+14x-8,
\]
and
\[
 \alpha=(1+t)(x-Y)-2,\qquad \beta=x-2,\qquad \gamma=-1.
\]
The preceding formulas give
\[
\begin{array}{lll}
q_{00}=2-12t,&q_{01}=6t(t+2),&q_{02}=-2(3t^2+3t-1),\\
q_{11}=-t(t^2+3t+3),&q_{12}=t(2t^2+6t+3),&
q_{22}=-t(t^2+3t+3).
\end{array}
\]
The division by the fixed line is exact already in $\QQ[t,x,Y]$:
\begin{align*}
 \Nm(s)-(Y-x)Q_s
  =-(Y+1)\bigl(Y^2+2Y-x^3+6x^2-14x+8\bigr).
\end{align*}
Substitution into \eqref{eq:appendix-determinant-expansion} gives
\begin{equation}
 4\det M_{1+t,1,1}=2t(24t^3+56t^2+33t+6).
\label{eq:appendix-simple-structural-factor}
\end{equation}
In particular, $4F_8(1,1,1)=12\ne0$ in $\QQ$.

\subsection*{A.2. The square-class identity}
Keep the normalized chart $e_1=v=1$ and put
\[
 r=u-1,\qquad d=\lambda-1,\qquad m=2c+e_2^3\lambda.
\]
Let $\widetilde F(z)=F_8(u,1,z)$, let
$\ell=\operatorname{lc}_z(\widetilde F)$, and set
$L_c=4\lambda^2\ell$.  Expanding
\eqref{eq:appendix-determinant-expansion} with the six coefficients in A.1
and collecting first in $d$ gives the following compact expression:
\begin{align}
L_c={}&3d^2e_3^2(3e_2^2r^2+3e_2r+1)\notag\\
&+3de_3(6e_2^2e_3r^2+3e_2e_3r-2e_2mr^2-3e_2mr-mr-2m)\notag\\
&+9e_2^2e_3^2r^2-6e_2e_3mr^2-9e_2e_3mr
  +m^2r^2+3m^2r+3m^2.
\label{eq:appendix-Lc-compact}
\end{align}
Define $N(z)=N_2z^2+N_1z+N_0$ by
\begin{align*}
N_2={}&-e_2(d+1)(6de_2e_3r+3de_3+6e_2e_3r-2mr-3m),\\
N_1={}&-(d+1)(3de_2^2e_3r^2-3de_2e_3r-3de_3
 +3e_2^2e_3r^2-3e_2e_3r+mr+3m),\\
N_0={}&r(d+1)(3de_2e_3r+2de_3+3e_2e_3r-m).
\end{align*}
Finally put
\[
 R(z)=2L_cz^3+R_2z^2+R_1z+R_0,
\]
where
\begin{align*}
R_2=3e_2\bigl[{}&d^2e_3(6e_2e_3r^3+3e_2+3e_3r^2-1)\\
&+d(12e_2e_3^2r^3+3e_2e_3-3e_2m+3e_3^2r^2
      -2e_3mr^3-3e_3mr^2+m)\\
&+6e_2e_3^2r^3-3e_2m-2e_3mr^3-3e_3mr^2\bigr],
\end{align*}
\begin{align*}
R_1=3\bigl[{}&d^2e_3(6e_2^3r^3+9e_2^3r^2+3e_2^2e_3r^4
 +3e_2^2r^2+15e_2^2r\\
&\hspace{5.3em}-3e_2e_3r^3-3e_2r+3e_2-3e_3r^2-1)\\
&+d(12e_2^3e_3r^3+18e_2^3e_3r^2+6e_2^2e_3^2r^4
 +3e_2^2e_3r^2+24e_2^2e_3r\\
&\hspace{3em}-2e_2^2mr^3-6e_2^2mr^2-6e_2^2mr
 -6e_2e_3^2r^3-3e_2e_3r+3e_2e_3\\
&\hspace{3em}-3e_2mr-3e_2m-3e_3^2r^2+e_3mr^3+3e_3mr^2+mr+m)\\
&+6e_2^3e_3r^3+9e_2^3e_3r^2+3e_2^2e_3^2r^4\\
&\hspace{3em}+9e_2^2e_3r-2e_2^2mr^3-6e_2^2mr^2-6e_2^2mr\\
&\hspace{3em}-3e_2e_3^2r^3-3e_2mr-3e_2m+e_3mr^3+3e_3mr^2\bigr],
\end{align*}
and
\begin{align*}
R_0={}&-3d^2e_3(6e_2^2r^3+3e_2^2r^2+3e_2e_3r^4
 +7e_2r^2+6e_2r+2e_3r^3+2r+2)\\
&+d(-36e_2^2e_3r^3-18e_2^2e_3r^2-18e_2e_3^2r^4
 -21e_2e_3r^2-18e_2e_3r\\
&\hspace{3em}+6e_2mr^3+12e_2mr^2+9e_2mr-6e_3^2r^3\\
&\hspace{3em}+3e_3mr^3+4mr^2+9mr+6m)\\
&-3r(6e_2^2e_3r^2+3e_2^2e_3r\\
&\hspace{3em}+3e_2e_3^2r^3-2e_2mr^2-4e_2mr-3e_2m-e_3mr^2).
\end{align*}
Let
\[
 \Delta=-4\left(q_{00}q_{11}-(q_{01}/2)^2\right),
 \qquad
 \widetilde F(z)=\sum_{k=0}^6 f_kz^k.
\]
In the ring
$\mathbf Z[c,\lambda,e_2,e_3,u,\lambda^{-1}]$, equating the seven
coefficients of $z$ in \eqref{eq:appendix-determinant-expansion} gives
\begin{align*}
16\lambda^2L_cf_6&=4L_c^2,\\
16\lambda^2L_cf_5&=4L_cR_2,\\
16\lambda^2L_cf_4&=R_2^2+4L_cR_1+3\Delta N_2^2,\\
16\lambda^2L_cf_3&=2R_1R_2+4L_cR_0+6\Delta N_1N_2,\\
16\lambda^2L_cf_2&=R_1^2+2R_0R_2+3\Delta(N_1^2+2N_0N_2),\\
16\lambda^2L_cf_1&=2R_0R_1+6\Delta N_0N_1,\\
16\lambda^2L_cf_0&=R_0^2+3\Delta N_0^2.
\end{align*}
These seven equalities are exactly the coefficients of $z^6,\ldots,z^0$
in the identity
\begin{equation}
 16\lambda^2L_c\,\widetilde F(z)=R(z)^2+3\Delta N(z)^2.
\label{eq:appendix-primitive-square-identity}
\end{equation}
They constitute the complete verification of
\eqref{eq:appendix-primitive-square-identity} in the stated polynomial ring.

Since $\ell=L_c/(4\lambda^2)$, define
\[
 A_{\mathrm c}=\frac{R}{2L_c},\qquad
 B_{\mathrm c}=\frac{N}{2L_c}.
\]
Dividing \eqref{eq:appendix-primitive-square-identity} by $4L_c^2$ gives
\begin{equation}
 \ell^{-1}\widetilde F=A_{\mathrm c}^2+3\Delta B_{\mathrm c}^2,
\label{eq:appendix-square-class-check}
\end{equation}
which is exactly \eqref{eq:norm-square-identity}.  Notice also directly from
the displayed leading terms that $\deg A_{\mathrm c}=3$ and
$\deg B_{\mathrm c}\le2$.

\subsection*{A.3. Geometric integrality of the residual octic}
For this argument, specialize to characteristic $13$ at
$(c,\lambda,e_1,e_2,e_3)=(1,1,1,0,2)$ and put
$H(u,v)=4F_8(u,v,1)$.  Expanding
\eqref{eq:appendix-determinant-expansion} gives
\begin{align*}
H(u,v)={}&6u^3v^5-u^3v^4+6u^3v^3-u^3v^2-u^3v
 -5u^2v^6+5u^2v^4+3u^2v^3+4u^2\\
&+5uv^7+3uv^6-3uv^4+3uv^3+4uv\\
&-6v^8-2v^7+2v^6+v^5-2v^4+4v^2.
\end{align*}
Its leading coefficient in $v$ is $-6=7\in\mathbf F_{13}^\times$, and
\[
 H(4,v)=7h(v),
 \qquad
 h(v)=v^8-3v^7-2v^6+3v^5+4v^4+4v^3-3v^2-5v-2.
\]
Irreducibility follows from the standard factor-degree criterion.  In
$\mathbf F_{13}[v]/(h)$ all powers are reduced by the single relation
\begin{equation}
 v^8=3v^7+2v^6-3v^5-4v^4-4v^3+3v^2+5v+2.
\label{eq:appendix-h-reduction}
\end{equation}
Put $r_j\equiv v^{13^j}\pmod h$, taking the representative of degree
at most seven.  Repeated use of \eqref{eq:appendix-h-reduction} gives
\begin{align*}
r_0={}&v,\\
r_1={}&-v^7+2v^6+6v^5-v^4+6v^3-2v^2-2v-6,\\
r_2={}&-4v^7+v^6-5v^5-3v^4+3v^3-3v^2-4v-1,\\
r_3={}&-3v^7-3v^6+4v^5+6v^3+2v^2+5v-2,\\
r_4={}&v^7+4v^6+2v^5+5v^4-4v^3-6v^2-2v+4,\\
r_5={}&-2v^7+4v^6+3v^5-4v^4-3v^3+4v^2-3v-4,\\
r_6={}&2v^7+5v^6-v^4-5v^3+4v^2+4,\\
r_7={}&-6v^7+3v^5+4v^4-3v^3+v^2+5v-5,\\
r_8={}&v.
\end{align*}
Thus $h$ divides $v^{13^8}-v$.  Also $r_4-v=r(v)$, where
\[
 r(v)=v^7+4v^6+2v^5+5v^4-4v^3-6v^2-3v+4.
\]
Moreover,
\[
 a(v)h(v)+b(v)r(v)=1
\]
with
\begin{align*}
a(v)&=5v^6+3v^5-3v^4-v^3-v^2-6v-2,\\
b(v)&=-5v^7+6v^6-5v^5+3v^4+2v^3-4v^2-2v-4.
\end{align*}
Hence $\gcd(h,v^{13^4}-v)=1$.  Recall that
$X^{q^n}-X$ is the product, without repetition, of the monic irreducible
polynomials over $\mathbf F_q$ whose degrees divide $n$.  Since $r_8=v$,
every irreducible factor of $h$ has degree dividing $8$; the displayed
B\'ezout identity excludes degrees dividing $4$.  The only divisors of $8$
are $1,2,4,8$, so every irreducible factor of $h$ has degree $8$.  Because
$\deg h=8$, the polynomial $h$ is irreducible.

If $H$ factored, the constant nonzero leading $v$-coefficient would force
each positive-$v$-degree factor to retain its $v$-degree after $u=4$, which
would factor $h$; a $v$-degree-zero factor would be a unit.  Thus $H$ is
irreducible.  Its homogenization is the specialized form $4F_8$.  Because
$H$ contains a nonzero $v^8$ term, $z$ does not divide that homogenization;
therefore any nontrivial homogeneous factorization would dehomogenize to a
nontrivial factorization of $H$.  The projective special fibre is thus
irreducible.  Finally
\[
 H(1,6)=0,
 \qquad
 (\partial_uH,\partial_vH)(1,6)=(2,0).
\]
Thus it has a smooth $\mathbf F_{13}$-point.  The field $\mathbf F_{13}$ is
perfect, so this irreducible curve is geometrically reduced.  If it were not
geometrically irreducible, Frobenius would act transitively on its geometric
components; a rational point would then lie on every component and would be
singular.  Hence the special octic is geometrically integral.  This is the
finite-field input in \cref{lem:residual-octic-integral}.

\subsection*{A.4. An admissible specialization}
For \cref{lem:admissible-nonempty} use
\[
 (c,\lambda,e_1,e_2,e_3)=(1,1,1,0,2),\qquad (u:v:z)=(3:5:1)
\]
over $\mathbf F_{13}$.  The shifted elliptic equation is
\[
 Y^2+2Y=x^3-6x^2+x+5,
\]
and the branch cubic $Y=0$ has discriminant $4$.  The fixed divisor is cut
out by
\[
 p(w)=w^3-w^2-2,
 \qquad \operatorname{disc}(p)=1,
\]
so it is reduced.  The residual conic has coefficient vector
\[
 (q_{00},q_{01},q_{02},q_{11},q_{12},q_{22})=(5,1,7,7,0,2)
\]
and factors as
\[
 5+x+7Y+7x^2+2Y^2
 =5(1+9x+Y)(1+12x+3Y).
\]
Thus the three norm lines are
\[
 L_0:Y=x,\qquad L_1:Y=4x-1,\qquad L_2:Y=4-4x.
\]
Substitution into the elliptic equation gives the three line-section cubics
\[
\begin{array}{c|c|c}
&\text{cubic in }x&\text{discriminant}\\ \hline
L_0&x^3+6x^2-x+5&4\\
L_1&x^3+4x^2+x+6&3\\
L_2&x^3+4x^2+2x-6&5.
\end{array}
\]
Hence all three intersections with $E$ are transverse.  To keep the sign
convention fixed, set
\[
 \mathcal R(x,Y)=x^3-6x^2+x+5-(Y^2+2Y),
\]
so that $E=V(\mathcal R)$.  The residual lines meet the branch line at
\[
 L_1\cap\{Y=0\}=(10,0),\qquad
 L_2\cap\{Y=0\}=(1,0),
\]
and
\[
 \mathcal R(10,0)=12,\qquad \mathcal R(1,0)=1.
\]
Their pairwise intersections are
\[
 L_1\cap L_2=(12,8),\qquad
 L_0\cap L_1=(9,9),\qquad
 L_0\cap L_2=(6,6),
\]
with
\[
 \mathcal R(12,8)=8,\qquad
 \mathcal R(9,9)=2,\qquad
 \mathcal R(6,6)=2.
\]
All five values are nonzero in $\mathbf F_{13}$.  Thus the residual line
sections avoid the branch divisor, and the three degree-three line sections
are pairwise disjoint on $E$.

The two residual line sections are therefore finite \emph{\'etale} away from
the branch locus and occur with multiplicity one in the norm divisor.  The
local product formula for the norm then forces exactly one of the three
sheets of $C\to E$ to carry a simple zero above each residual point, which is
condition~(iii) of \cref{def:admissible-locus}.

The canonical-model condition follows from the same functions used in
\cref{prop:theta-resolution}.  Because $\mathbf F_{13}$ contains $\zeta_3$,
the nine quadratic
representatives split into
\[
 \langle1,x,Y,x^2\rangle,\qquad
 \langle w,xw,wY\rangle,\qquad
 \langle w^2,xw^2\rangle.
\]
Their pole orders at a point over $O$ are $(0,2,3,4)$, $(1,3,4)$, and
$(2,4)$, respectively.  Thus the only canonical quadric is
$X_0X_2-X_1^2=0$, of rank three, and $|\kappa|$ is the unique $g^1_3$.
A member of this pencil is cut out by $a+bw$, so
$\Nm_\pi(a+bw)=a^3+b^3Y$; in particular its norm line has zero
$x$-coefficient.  By contrast, $L_0,L_1,L_2$ all have nonzero
$x$-coefficient.  Their inverse degree-three
divisors are therefore not linearly equivalent to $\kappa$, hence lie in the
$h^0=1$ Abel locus and map to the smooth theta locus.  This verifies
condition~(iv) and completes the arithmetic specialization used in
\cref{lem:admissible-nonempty}.

\begin{thebibliography}{99}

\bibitem{BBD1982}
A.~A. Beilinson, J. Bernstein, and P. Deligne,
\emph{Faisceaux pervers}, Ast\'erisque \textbf{100}, Soci\'et\'e Math\'ematique de France, Paris, 1982.

\bibitem{BirkenhakeLange2004}
Christina Birkenhake and Herbert Lange,
\emph{Complex abelian varieties}, second ed., Grundlehren der mathematischen Wissenschaften, vol.~302, Springer-Verlag, Berlin, 2004.

\bibitem{BrasseletLeSeade2000}
Jean-Paul Brasselet, D\~ung Tr\'ang L\^e, and Jos\'e Seade,
\emph{Euler obstruction and indices of vector fields}, Topology \textbf{39} (2000), no.~6, 1193--1208.

\bibitem{deCataldoMigliorini2002}
Mark Andrea A. de Cataldo and Luca Migliorini,
\emph{The hard Lefschetz theorem and the topology of semismall maps}, Ann. Sci. \`Ecole Norm. Sup. (4) \textbf{35} (2002), no.~5, 759--772.

\bibitem{Dynkin1957}
E.~B. Dynkin,
\emph{Semisimple subalgebras of semisimple Lie algebras}, Amer. Math. Soc. Transl. Ser.~2 \textbf{6} (1957), 111--244.

\bibitem{FraneckiKapranov2000}
J. Franecki and M. Kapranov,
\emph{The Gauss map and a noncompact Riemann--Roch formula for constructible sheaves on semiabelian varieties}, Duke Math. J. \textbf{104} (2000), no.~1, 171--180.

\bibitem{FultonHarris1991}
William Fulton and Joe Harris,
\emph{Representation theory: A first course}, Graduate Texts in Mathematics, vol.~129, Springer-Verlag, New York, 1991.

\bibitem{EGA4}
Alexander Grothendieck and Jean Dieudonn\'e,
\emph{\'El\'ements de g\'eom\'etrie alg\'ebrique. IV. \'Etude locale des sch\'emas et des morphismes de sch\'emas, troisi\`eme partie}, Publ. Math. Inst. Hautes \'Etudes Sci. \textbf{28} (1966), 5--255.

\bibitem{Hartshorne1977}
Robin Hartshorne,
\emph{Algebraic geometry}, Graduate Texts in Mathematics, vol.~52, Springer-Verlag, New York, 1977.

\bibitem{JKLM2025}
Ariyan Javanpeykar, Thomas Kr\"amer, Christian Lehn, and Marco Maculan,
\emph{The monodromy of families of subvarieties on abelian varieties}, Duke Math. J. \textbf{174} (2025), no.~6, 1045--1149.

\bibitem{Kashiwara1985}
Masaki Kashiwara,
\emph{Index theorem for constructible sheaves}, in \emph{Syst\`emes diff\'erentiels et singularit\'es}, Ast\'erisque \textbf{130} (1985), 193--209.

\bibitem{Kempf1973}
George Kempf,
\emph{On the geometry of a theorem of Riemann}, Ann.\ of Math. (2) \textbf{98} (1973), no.~1, 178--185.

\bibitem{KLM2026}
Thomas Kr\"amer, Christian Lehn, and Marco Maculan,
\emph{Exceptional Tannaka groups only arise from cubic threefolds}, J. Reine Angew. Math. \textbf{830} (2026), 51--99.

\bibitem{KraemerWeissauer2015}
Thomas Kr\"amer and Rainer Weissauer,
\emph{Vanishing theorems for constructible sheaves on abelian varieties}, J. Algebraic Geom. \textbf{24} (2015), no.~3, 531--568.

\bibitem{Laufer1977}
Henry B. Laufer,
\emph{On minimally elliptic singularities}, Amer. J. Math. \textbf{99} (1977), no.~6, 1257--1295.

\bibitem{Litt13}
Daniel Litt,
\emph{Problem no.~13}, Problems I Like, \url{https://www.problemsilike.com/13}, accessed August 10, 2026.

\bibitem{Macdonald1962}
I.~G. Macdonald,
\emph{Symmetric products of an algebraic curve}, Topology \textbf{1} (1962), 319--343.

\bibitem{Pinkham1977}
Henry C. Pinkham,
\emph{Simple elliptic singularities, del Pezzo surfaces and Cremona transformations}, Several Complex Variables, Proc. Sympos. Pure Math., vol.~30, American Mathematical Society, Providence, RI, 1977, pp.~69--71.

\bibitem{Saito1988}
Morihiko Saito,
\emph{Modules de Hodge polarisables}, Publ. Res. Inst. Math. Sci. \textbf{24} (1988), no.~6, 849--995.

\bibitem{stacks-project}
The Stacks Project Authors,
\emph{Stacks Project}, \url{https://stacks.math.columbia.edu}, accessed August 10, 2026.

\end{thebibliography}

\end{document}


